% SOURCE docs/release/2026-10-29/source/docs/writing-repair/2026-10-29/src/ef4782d6b51eaaa8.tex CHAR 0
\chapter{Choosing tools for an accountable research assistant}
\label{exec14:l6}
\label{ch:la05-prototype}

The workflow comparison becomes a component decision here. The opening
frameworks and historical local controller establish the initial boundary.
The subsequent studies test evaluation, memory, information exchange, code
assistance and scientific reasoning against that boundary. Broader framework
and backend profiles then lead into their detailed parsing, validation, retry
and recovery paths. The question throughout is which useful operation can
be admitted without changing the causal object or overstating its evidence.


\section{Define the research task before choosing the software}

The proposed causal assistant has three separable responsibilities. It helps
a person express a banking decision, invokes mathematical operations whose
applicability is explicit, and preserves the evidence needed to interpret the
answer. A model backend proposes text or structured objects. An orchestration
framework arranges calls and state transitions. A causal engine supplies a
particular graphical or numerical calculation. Their outputs have different
evidential meanings even when one software package exposes all three through
the same chat interface. The architecture should make those meanings visible
before selecting a vendor or adopting a framework.

Consider a lender comparing an offer with no offer for accounts entering
arrears. The declared population, twelve-month outcome, information available
at first eligibility, treatment contrast and cost convention determine the
question. A model may propose a graph, explain why a field is measured too
late or draft an analysis plan. A numerical service may compute a conditional
effect under the supplied graph and sampling law. The controller must prevent
the first activity from silently certifying the assumptions of the second.
The servicing comparison in Section~\ref{exec13:agent-architecture} develops
the substantive causal alternatives; here the question is which software
boundary preserves them.

Let an analysis state be
\begin{equation}
 s=(q,d,g,a,m,r,e),
 \label{exec14:l6-state}
\end{equation}
where $q$ is the decision contract, $d$ the data and observation definitions,
$g$ the graph or graph class, $a$ the assumptions and their statuses, $m$ the
method request, $r$ the result and $e$ its provenance. Each component has a
version identity. An external grant $u$ specifies permitted operations and
inputs. It is not a field that the model can authorize for itself. A proposed
transition $f$ should be admitted only if its preconditions hold for the
current state and grant. The output must retain the identities it actually
used. These are proposed application invariants, not properties supplied by
the word agent in a package description. \citep{exec14_l6_design}

This formulation also prevents an unhelpful choice between unrestricted
automation and manual work. Some transitions are mechanical: parsing a
versioned request, verifying source bytes or detecting a graph cycle. Others
need scientific judgment: whether an omitted variable is plausibly a common
cause, whether a policy's information set reflects actual availability, or
whether a transport assumption is defensible. An admissible workflow can
automate the first group and record uncertainty in the second. A fixed
controller can continue independent analyses after one method refuses a
request, provided it does not substitute a new target without an explicit
revision.

\section{Backend alternatives and what can actually be compared}

An offline fixture is the first useful backend. It returns a predetermined
proposal whose intended errors and ambiguities are known. It can establish
whether the controller preserves a PAG, refuses an unsupported policy and
binds a result to the right horizon. It cannot establish how often a language
model generates a correct proposal. The existing local prototype accepts
authored or recorded proposals and makes no live model calls. Recorded
provider/model labels are caller-supplied provenance, not independently
authenticated proof of generation. \citep{exec13_architecture_local,exec14_l6_local}

A prospective local model moves the inference computation to a selected
host. Qualification would bind weight files, tokenizer, inference software,
quantization and decoding settings, model terms and any remotely retrieved
material. A prospective hosted model binds a provider endpoint, model
revision or declared alias policy, request/response schema, data transfer
boundary, usage accounting and service contract. These are comparison
requirements. Neither local execution nor a commercial contract by itself
establishes causal competence, confidentiality of every component, or
permission to use a particular source. Current model rankings, prices and
retention guarantees are not inferred from generic framework examples.
\citep{exec14_l6_design}

\begin{longtable}{@{}P{.19\textwidth}P{.34\textwidth}P{.40\textwidth}@{}}
\caption{Backend choices for the proposed assistant. Except for the offline
controller evidence, these are prospective contracts, not measured services.}
\label{exec14:l6-backends}\\
\toprule Choice & Main benefit to investigate & Evidence needed before use\\\midrule
\endfirsthead
\toprule Choice & Main benefit to investigate & Evidence needed before use\\\midrule
\endhead
Authored or recorded proposal & Deterministic fault and semantic cases without
provider cost & Exact bytes, declared origin and an independent expected
disposition; no model-quality inference.\\
Local model & A controlled inference environment and inspectable deployment
inputs & Weight/tokenizer terms and hashes, runtime dependencies, outbound
access and logging, capacity, task performance and full operating cost.\\
Hosted model & A service interface with outsourced inference operations &
Exact service/model identity, allowed inputs, supported schema behavior,
usage and price evidence, retention and recovery contract.\\
Selective routing & Use additional model work only for eligible hard cases &
A frozen routing rule, measured conditional error and cost, comparable
information and an explicit unavailable route when no eligible backend exists.\\
Human or deterministic baseline & An interpretable reference for the value
of assistance & The same task definitions, tools and total repair opportunity;
record review time and failure dispositions.\\
\bottomrule
\end{longtable}

The common envelope should preserve the exact input and response bytes,
provider or local runtime identity, relevant settings, tool definitions,
retrieval identities, schema version and measured usage when available.
An adapter should distinguish malformed structure, an explicit backend
refusal, truncation, timeout and an unknown remote outcome. It should not
replace any of these with an empty successful causal result. A structurally
valid graph can still contain unsupported causal edges, so structural
validation and scientific adjudication remain separate.

Model substitution is therefore a reviewed change to the analysis system.
If a provider alias or local weight file changes, previously saved outputs
remain evidence of their original run. A new run should receive a new
execution identity, even if it preserves the same business question.
Conversely, keeping the same model while changing the target population
requires a new scientific identity. Requiring these two distinctions avoids
both needless recomputation of unrelated artifacts and unsafe reuse of
results for a changed question. The custody construction in
Section~\ref{exec14:artifact-custody} describes the narrower local byte and
context checks already observed.

\section{Three orchestration frameworks under inspected revisions}

LangGraph offers a state-graph and checkpoint interface. At the inspected
revision dated 7 October 2026, its project declares version 1.2.14 and Python
3.10 or later. The selected compile path accepts a checkpointer. Its interrupt
documentation is precise about resumption: the interrupted node starts again
and re-executes its preceding logic. That behavior is useful for a review
step, but an expensive or externally consequential action before the
interrupt may be repeated unless the application separately journals or
deduplicates it. The proposed reuse decision is a conditional Python
adapter for state orchestration, with scientific identities and side-effect
records owned by the application. Porting the entire framework is not
justified by the present proposal. The inspected root license is MIT;
dependency and deployment review remains separate. \citep{v12_la04_langgraph}

Pydantic AI offers typed tool interfaces and usage controls. The inspected
7 October 2026 revision declares Python 3.11 or later and a dynamic package
version; no exact release number is inferred from that field. Selected tool
fields support sequential execution and an approval condition, but their
inspected defaults do not require either. The application must configure
consequential tools explicitly. Request counts can be checked before a model
call, whereas token totals generally arrive with responses; an optional
pre-counting path has its own support and overhead. The selected cost check
warns when no cost was calculated and cannot enforce the stated monetary
limit in that case. The proposed decision is conditional reuse for typed
adapters, backed by an application-owned budget ledger and unknown-cost
policy. The MIT notice does not resolve every model or data term.
\citep{v12_la04_pydantic_ai}

AutoGen's inspected April 2026 snapshot declares AgentChat 0.7.5 and a pinned
core dependency. Its README describes maintenance mode with no new features;
that is a source declaration, not an independently verified evaluation of
any successor product. The selected AssistantAgent is stateful and expressly
not safe for concurrent calls on the same instance. It can execute returned
tool calls concurrently, and reflection can introduce another inference
call. Its save/load methods preserve model context; they do not record the
outcome of every external action. The code notice inspected for AgentChat
is MIT while the root documentation notice is CC-BY-4.0. The proposal retains
AutoGen as a comparative architecture and possible isolated adapter, but
does not make it the default controller or treat conversation persistence
as an exactly-once journal. \citep{v12_la04_autogen}

\begin{longtable}{@{}P{.20\textwidth}P{.35\textwidth}P{.38\textwidth}@{}}
\caption{Framework decisions concern selected source paths; none of these
frameworks has been executed or qualified by this comparison.}\\
\toprule Candidate & Useful inspected mechanism & Condition on adoption\\\midrule
\endfirsthead
\toprule Candidate & Useful inspected mechanism & Condition on adoption\\\midrule
\endhead
LangGraph & State graph, checkpoint and interrupt & Preserve immutable
scientific inputs; make replay of pre-interrupt work deliberate.\\
Pydantic AI & Tool schemas, control fields and usage checks & Configure
authorization/sequencing; reserve shared budgets; refuse unsupported unknown
cost when a hard ceiling is required.\\
AutoGen & Stateful model context and tool/reflection loop & Isolate mutable
instances; account for concurrency and reflection; review maintenance and
code/documentation terms.\\
Plain local controller & Fixed supported operations and a narrow journal &
Retain the observed limited scope; add capabilities only with their own
scientific and operating contracts.\\
\bottomrule
\end{longtable}

The plain controller is the baseline because its admitted operations and
failure states can be enumerated. This is not evidence that frameworks add
no value. A framework may reduce implementation effort for complex state,
streaming or human review. That advantage should be evaluated against the
additional dependencies, semantics and operational responsibilities it
introduces. A thin adapter can preserve the application's causal request
and result contracts while allowing this comparison. No framework should
be credited with discovery, identification or calibrated uncertainty merely
because it can call a tool bearing those names.

\section{Causal packages need decisions at the scientific boundary}

The package inventory covers different kinds of reusable material. PyWhyLLM
is a hypothesis-elicitation candidate: its inspected back-door suggestion
path delegates to a model, while the identification-suggester front-door
path is a placeholder. A separate simple suggester can return front-door
names; neither path supplies a formal identification certificate, as
Section~\ref{exec14:llm-identification} explains. A suggested confounder is
not a proven adjustment set. Adopt
the idea of source-linked suggestions conditionally, and pass a supplied
graph to a qualified identification engine. Directly porting the entire
dependency stack would add maintenance without resolving that distinction.
\citep{v12_la_package_pywhy_llm}

The inspected CausalAgent package supplies interaction and tool-wrapper
ideas, but its OLC handler's missing function binding and permissive
malformed-output path prevent default direct adoption of that route. A
wrapper must carry typed errors and graph endpoint semantics. In oci-agent,
the useful intermediate objects are specifications, notebooks and diagnostic
reports. Its inspected unrestricted notebook runner has no finite per-cell
timeout, and its selected trimmed AIPW target differs from conventional
overlap weighting. Adapt the specification/report concepts behind bounded
operations; do not expose the general runner as a trusted causal endpoint.
These are static source decisions, not measured failure frequencies.
\citep{v12_la_package_causalagent,v12_la_package_oci_agent}

Benchmark software has another role. CLadder and CausalDS are candidates
for evaluating task formulations and answers after their generator/scorer
assumptions are reviewed. They are not engines whose test labels establish
the bank's causal truth. The selected IV-as-ATE and unknown-identification
boundaries remain reasons to quarantine unsupported cases or repair a
scorer before using it. OntoLearner can contribute structured vocabulary
and local versioned ontology objects, subject also to the document-identity
checks in Chapter~\ref{exec17:ontology}; taxonomy does not certify a causal
edge. LCM remains a deferred execution candidate until weights,
terms and its sufficient lagged regime are qualified. The LLM-CD and
LLM-SCD code pathways retain unresolved license or version-alignment
requirements; their papers can still be evaluated as scientific arguments.
\citep{v12_la_package_cladder,v12_la_package_causalds,v12_la_package_ontolearner,
v12_la_package_lcm,v12_la_package_llm_cd,v12_la_package_llm_scd}

These choices distinguish five useful dispositions. Direct adoption requires
the actual component's contract, environment and rights to match the use.
Conditional integration adds an adapter and predeclared acceptance cases.
Porting reimplements selected semantics and needs new numerical and error
comparisons; source similarity is insufficient. Reference-only use borrows
an idea without claiming executable reuse. Rejection or deferral records
the reason and a concrete reconsideration trigger. A package can receive
different dispositions for its graph representation, estimator and notebook
runner. Treating its repository as a single yes-or-no choice obscures the
component that creates the real scientific or operating risk.

\section{A state invariant for revision and invalidation}

Suppose the controller maintains a dependency graph of immutable analysis
objects. A result stores the exact identities of the target, input data,
graph, method and execution environment used to produce it. Define the
reuse invariant $I(s)$ to mean that every result presented as current has
matching admitted dependencies and retains its qualification status.
Initially, no result exists and $I$ holds. A transition that creates a
result establishes $I$ only after checking its dependencies. A transition
that changes a dependency preserves $I$ by retiring affected descendants
from the current view and requiring recomputation or an explicit stale
label. By induction, those transition rules preserve $I$ for a finite
sequence. This is an application-level argument under the stated rules,
not a proof that the current prototype implements an arbitrary dependency
graph or that supplied causal assumptions are true. \citep{exec14_l6_exact}

The retirement rule should be selective. In a finite illustrative dependency
graph, identification depends on the target and graph; estimation also
depends on data; valuation depends on the estimate and cost assumptions;
the report depends on valuation. A pure observation-clock diagnostic depends
on the data and clock rule. Changing the target horizon invalidates
identification, estimation, valuation and the report, while leaving that
diagnostic current if its own dependencies did not change. Changing costs
invalidates valuation and the report without changing the effect estimate.
An exact traversal verifies these sets in the project example. No claim is
made that every real diagnostic is target-independent: if its admissibility
rule uses the horizon, that dependency must be recorded too.
\citep{exec14_l6_exact}

This rule protects a banking decision from a subtle error. An agent can
produce a plausible revision to the question after observing a weak result.
The old estimate may remain numerically unchanged and its source files may
still pass all hashes. It becomes ineligible as an answer to the new
question because its scientific dependencies changed. The revised target
and the original result should both remain inspectable. The controller
must also prevent a later successful parse, retry or report step from
turning a conditional or unverified premise into a verified one. Such
promotion requires new supporting evidence and a separate recorded decision.

\section{Recovery includes an unknown outcome}

Local pure computations and external effects have different recovery
contracts. A deterministic calculation can be rerun with pinned inputs and
environment when recomputation is allowed. A paid model request may have
completed and been billed even if its response was lost. A proposed journal
distinguishes prepared, submitted, acknowledged, completed, failed and
outcome-unknown states. A timeout after submission should not be treated as
proof that the provider did no work. The availability of provider lookup or
idempotency must be established for the selected service.
\citep{exec14_l6_design}

There is a simple indistinguishability argument. Let history $h_0$ lose a
request before the external service receives it, and let history $h_1$
lose the response after the service completes the work. If the local
journal in each case contains only submitted and the same timeout, its
observable state is identical. A rule depending only on that state takes
the same action in both histories. Resubmission can duplicate work under
$h_1$; refusing to resubmit can leave the task unfinished under $h_0$.
Randomization changes the probabilities but cannot guarantee both no
duplication and eventual completion from this information alone. A service
lookup, idempotent external operation or an explicitly accepted duplication
policy adds information or changes the requirement. This finite argument
concerns the stated observation model, not all possible distributed
systems. \citep{exec14_l6_exact}

For a causal assistant, the appropriate response may be to retain an
unknown provider outcome while continuing independent deterministic checks
or another task. Cancellation similarly needs a scope: stopping future
local work does not prove that an already submitted remote request was
cancelled or its charge reversed. The budget ledger should retain any
reservation whose financial outcome is unknown until reconciliation. The
scientific report can remain incomplete with an explicit reason; filling
the missing answer from memory would sever its evidence chain.

\section{Shared budgets require reservations}

Usage checks on individual calls do not by themselves enforce a shared
budget under concurrency. Let $B$ be a stipulated expenditure ceiling,
$S$ settled cost and $R$ the sum of unresolved reservations. Before admitting
request $j$ with a credible upper cost bound $\bar c_j$, a controller can
require
\begin{equation}
 S+R+\bar c_j\le B.
 \label{exec14:l6-reserve}
\end{equation}
The check and reservation must be one serialized state transition. If the
actual cost $c_j$ is no larger than its reserved bound, completion replaces
that reservation with $c_j$; hence the total does not increase beyond its
admitted level. Starting from $S=R=0$, induction preserves $S+R\le B$.
Unknown prices, unbounded usage, excluded fees or provider overcharges
invalidate the premise; they are not resolved by this algebra.
\citep{exec14_l6_exact}

An exact illustrative case makes the race visible. Take $B=10$ dollars and
settled cost $S=6$. Two requests each require a three-dollar upper bound.
If each reads the same unreserved balance, both individually see
$6+3\le10$, yet their combined cost can reach 12 dollars. Under serialized
reservation, the first makes $S+R=9$, and the second is refused because
$6+3+3>10$. If the first settles at two dollars, $S$ becomes eight and $R$
returns to zero; the second three-dollar reservation is still inadmissible.
The project calculation enumerates both admission orders and verifies the
same conclusion. All amounts are stipulated; they are not prices of any
model. \citep{exec14_l6_exact}

A successful provider response may supply token counts after money has
already been spent. Accounting can record that cost correctly without
having prevented an overspend. Conversely, a token cap can limit one
resource while leaving human review time or a separately priced tool
unbounded. A useful prototype records call counts, known usage, reservations,
settled expense and unknown charges separately. Retries and reflection
consume the same shared budget, and failed requests remain in its total.
These proposed controls address the inspected framework boundaries without
claiming that a financial guarantee follows from a field named cost limit.

\section{Routing is a decision rule, not a model leaderboard}

A prospective router can compare eligible backends by expected task loss
and full incremental cost. For task information $z$, let $e_b(z)$ be a
calibrated error probability for backend $b$, $\ell(z)$ the loss of the
defined error, and $c_b(z)$ its expected call and review cost. Under the
deliberately simplified assumption of one relevant error category and a
fixed loss for that category, choose among eligible backends by
\begin{equation}
 b^*(z)\in\arg\min_{b\in\mathcal B(z)}
       \{c_b(z)+\ell(z)e_b(z)\}.
 \label{exec14:l6-route}
\end{equation}
The eligible set $\mathcal B(z)$ encodes data-use and capability constraints.
An ineligible backend cannot become eligible by having a low quoted cost.
If the set is empty, the workflow returns an unavailable route or uses an
authorized non-model alternative. Real error categories, uncertain losses
and abstention require a richer decision model. \citep{exec14_l6_exact}

For illustration, backend A costs one dollar with error probability .08,
and backend B costs four dollars with error probability .02. Then B has
lower expected total cost precisely when
$4+.02\ell<1+.08\ell$, or $\ell>50$ dollars. At loss 100, the stipulated
totals are nine and six dollars; at loss ten they are 1.80 and 4.20 dollars.
These numbers explain what a pilot should measure. They neither rank real
models nor justify applying a benchmark-wide error rate to a new banking
population. Costs of adjudication and wrong refusals must be included if
they differ between routes. \citep{exec14_l6_exact}

To evaluate a router, freeze its task strata and eligibility rules before
evaluation, give alternatives comparable information and tools, and retain
every assigned case. Selecting the backend after seeing the correct answer
would evaluate an oracle unavailable at deployment. A simple deterministic
form, a single assisted controller and a critic-assisted controller provide
useful alternatives before sophisticated routing. The experiments should
measure whether the extra model work changes target errors, unresolved
assumptions, usable results and review effort. Scientific identification
still belongs to the causal model and data; routing changes assistance,
not what the observational law identifies.

\section{From architectural comparison to a bounded prototype}

The proposed adoption sequence begins with the existing fixed-operation
controller and generated inputs. A selected framework can then be placed
behind an adapter preserving the same requests, grants, result types and
failure meanings. Offline fixtures test schema and recovery behavior before
an actual model is admitted. A live trial, if justified, needs named model
and service identities, permitted data, measured usage and an independent
task rubric. It can use public and generated tasks to establish whether
further bank experimentation is worth funding. No authentic customer
trial is required merely to compare the proposed interfaces.

This opening comparison establishes an initial scope through three pinned
orchestration frameworks and selected causal packages. The preserved local
implementation that follows provides a baseline. The later framework and
backend profiles in Sections~\ref{sec:exec24-framework-census} and
\ref{sec:exec24-backend-census}, together with their detailed client and
delegate studies, broaden the source comparison. They still do not establish
complete coverage or native qualification. Model terms, deployment rights,
task performance, costs and authentic owner acceptance remain distinct
questions; feature and price claims require their own applicable evidence.
\citep{exec14_l6_design,exec14_l6_local}

\label{exec14:l6-extension-end}


% SOURCE docs/release/2026-10-29/source/docs/writing-repair/2026-10-29/src/ff45440f911890d9.tex CHAR 0
\section{An adapter must preserve the scientific object}
\label{exec14:l6-adapter}

Typed input is useful only if the adapter preserves what the fields mean.
Consider two implementations of the same analysis path: a direct local
service and a framework-mediated service. Let $x$ be a valid request with
its scientific context, and let $F(x)$ be the result of a qualified
deterministic engine under that context. An adapter encodes the request
with $E$, sends it through the tool interface, and decodes the response
with $D$. The desired comparison is semantic:
\begin{equation}
 D\{F_{\rm tool}(E(x))\}\equiv F(x),
 \label{exec14:l6-commute}
\end{equation}
where $\equiv$ requires the same target, assumptions, qualification and
result meaning, not identical incidental ordering of JSON fields. For a
pure deterministic numerical path it may also require exact values, or a
declared numerical tolerance supported by the computation. For stochastic
methods the comparison must declare randomization and distributional
criteria. Passing one exact fixture does not prove this relation for every
request. This is a proposed qualification relation with explicit domains,
not an asserted property of an installed framework.
\citep{exec14_l6_adapter_exact}

Several superficially convenient encodings cannot satisfy that relation.
Suppose an encoder replaces every finite or unbounded effect set by a pair
containing its smallest and largest endpoints. Then the union
$(-\infty,-2]\cup[3,\infty)$ and the whole line both produce the same
pair of infinite endpoints. A downstream decoder receiving only that pair
cannot know which input was supplied. If it returns the union, it is wrong
for the whole-line input; if it returns the whole line, it is wrong for the
union as an exact representation. It could return a conservative outer
approximation if the contract explicitly permits loss of precision, but
it must label that transformation rather than claim equivalence. An exact
finite representation of endpoint pairs reproduces this collision in the
project example.

There is a general reason. Let $E(x_1)=E(x_2)$ for two scientifically
distinguishable objects $x_1$ and $x_2$. Any deterministic decoder $D$ has
the same output on the two equal encodings. Therefore both identities
$D(E(x_1))=x_1$ and $D(E(x_2))=x_2$ cannot hold. Randomized decoding cannot
guarantee recovery of both without additional information either. The
same argument applies when an adapter collapses PAG endpoint types, maps
unknown and nonidentified into one status, or drops which population an
effect targets. Before considering a port, the proposal should inspect
whether the receiving schema is expressive enough. More elaborate retry
logic cannot reconstruct information the schema discarded.
\citep{exec14_l6_adapter_exact}

A refusal state also carries scientific information. A graphical engine
can return an identification expression, a formal nonidentification result
within a specified model class, an unsupported method, or an operational
failure. These states call for different continuation actions. A framework
that treats every returned string as a successful tool result may preserve
transport bytes while discarding that distinction in its next prompt.
The adapter should expose a tagged result and require the next operation
to consume an admitted tag. A report can state that a method is unsupported
without claiming the target is unknowable. A timeout cannot become an
estimated effect of zero. These conditions can be tested with authored
fixtures before a model is selected.

\subsection{Units, signs and set-valued banking decisions}

Unit conversion illustrates a transformation that can preserve meaning
when it is explicit. Let an effect $\theta$ be measured in dollars and
suppose an interface returns thousands of dollars: $\theta'=\theta/1000$.
An interval or union must be transformed as a set, not only by changing
a display label. More generally, for a nonzero constant $a$ and constant
$b$, let $T(\theta)=a\theta+b$. For a closed interval $[l,u]$, its image
is $[al+b,au+b]$ when $a>0$ and $[au+b,al+b]$ when $a<0$. Infinite
endpoints and each component of a union follow the same ordering rule.
Empty and whole-line sets retain their distinct identities.
\citep{exec14_l6_adapter_exact}

If a set $C$ has coverage at least $1-\alpha$ for $\theta$, then its exact
image $T(C)$ has coverage at least $1-\alpha$ for $T(\theta)$ when $T$ is
a fixed known bijection: the events $\theta\in C$ and
$T(\theta)\in T(C)$ are identical. An outer enclosure of $T(C)$ can
preserve the lower coverage guarantee while losing precision. A numerical
inner approximation has no such automatic guarantee. If the conversion
rate or cost is estimated, the fixed-known-map premise fails and its
uncertainty must be accounted for. This elementary argument explains why
representation and numerical error belong in an adapter contract. It
does not validate the original set's sampling law.

Take a stipulated offer-minus-no-offer effect set
$C=(-\infty,-2]\cup[3,\infty)$ dollars and an incremental offer cost
of one dollar. The net-value set is
$C-1=(-\infty,-3]\cup[2,\infty)$. Its worst case is negative and its
best case is positive, so this set alone does not justify either robust
acceptance or robust rejection under a sign rule. Reporting only the
midpoint of a finite substitute interval would invent a decision. If a
receiving tool instead uses the reverse treatment contrast, the effect
set becomes $-C=(-\infty,-3]\cup[2,\infty)$, but this numerical
coincidence in the chosen example does not make reversing the contrast
equivalent to subtracting cost. The object must retain the target and
transformation, not only the endpoints.

For a less coincidental bounded check, let $C=[2,5]$ and a positive
one-dollar cost. Net value is $[1,4]$, so every admitted value supports
the offer under the stipulated expected-value sign criterion. Reversing
the effect contrast produces $[-5,-2]$; it is not $[1,4]$. Scaling
dollars into thousands maps the original effect to $[.002,.005]$ and
the cost to $.001$ in those units. Subtracting a cost of one in the new
units would be a thousandfold error. Exact rational images in the example
check the sign, units and component ordering. The resulting conclusions
remain illustrative consequences of stipulated sets and costs, not
validated lending advice or empirical financial benefit.
\citep{exec14_l6_adapter_exact}

\subsection{Porting needs a failure-domain comparison}

A port's acceptance domain should include both supported and deliberately
unsupported requests. In supported cases compare scientific context,
endpoint conventions, uncertainty tags, set topology, estimates and
applicable numerical tolerances. In unsupported cases compare the reason
and retained provenance. Test NaN, infinity, missing timestamps, duplicate
variable identifiers, a changed sign convention and a graph with unresolved
endpoints according to the receiving method's contract. The expected
answer may be refusal. A test that merely requires the port to return
something for every case rewards the very semantic substitutions the
product is meant to prevent.

The comparison should record which differences are representational and
which are substantive. Object-key order may be irrelevant to the engine
yet consequential to a byte digest; a normalized serialization can be
introduced with an explicit version and migration. Changing a dependency
law from independent accounts to monthly panels is substantive even if
the numerical array is unchanged. A port can preserve syntax and still
change inference. These distinctions motivate a thin typed adapter first,
followed by a port only when measured operational benefits justify the
additional qualification burden.

\section{Attribute improvements to the component that changed}
\label{exec14:l6-factorial}

Comparing an elaborate new agent to an old script can conflate at least
three changes: a stronger backend, better orchestration and access to
better causal tools. A product decision needs to know which component
accounts for an observed improvement and at what cost. A prospective
comparison can hold documents, permitted tools, review rubric and total
repair opportunity fixed while varying backend and orchestration. Let
$b\in\{A,B\}$ index two candidate backends and $f\in\{0,1\}$ a plain
controller versus the framework candidate. For eligible task $i$, define
loss $L_i(b,f)$ to include the predeclared scientific error and operational
failure consequences. Lower loss is better.

If all four outcomes can be obtained without contamination, a paired
task comparison isolates the framework contrast at each backend. If
repeated exposure changes human judgment or leaks a reference answer,
use independent adjudicators, blinded outputs or randomized assignment
with a declared estimand. Neither approach permits tuning the framework
on the final task answers. Tasks that share a generating mechanism are
not automatically independent observations, so task families and their
sampling law must be declared before attaching inferential uncertainty.
The prospective construction here is an evaluation design, not a claim
of collected experimental outcomes.
\citep{exec14_l6_adapter_exact}

Write $\mu_{bf}=E{L_i(b,f)}$ for the mean under the target evaluation
task distribution. The mean framework contrast, averaging equally over
the two backends, is
\begin{equation}
 \Delta_F=\tfrac12[(\mu_{A1}-\mu_{A0})+
                              (\mu_{B1}-\mu_{B0})],
\end{equation}
and the analogous backend contrast is
\begin{equation}
 \Delta_B=\tfrac12[(\mu_{B0}-\mu_{A0})+
                              (\mu_{B1}-\mu_{A1})].
\end{equation}
The interaction
$\Delta_{BF}=(\mu_{B1}-\mu_{B0})-(\mu_{A1}-\mu_{A0})$
measures how the framework contrast changes with backend. These are
descriptive contrasts of the stipulated factorial task experiment. Their
identification from sampled trials depends on the assignment and missing
outcome rules. Unequal deployment frequencies would require the relevant
weights rather than the equal weights used here.

An exact example takes mean losses
$(\mu_{A0},\mu_{A1},\mu_{B0},\mu_{B1})=(10,8,7,6)$.
Then $\Delta_F=-1.5$, $\Delta_B=-2.5$ and $\Delta_{BF}=1$.
The framework reduces mean loss by two units with A but only one with B.
The diagonal comparison $\mu_{B1}-\mu_{A0}=-4$ mixes both changes;
it cannot all be attributed to orchestration. The average contrasts
sum to the diagonal difference, yet their separate magnitudes carry the
investment information. Costs require their own four-cell measurement,
including failed tasks and review effort. The numbers are authored
illustrations, not performance claims about any named framework.
\citep{exec14_l6_adapter_exact}

Availability creates a further denominator problem. Suppose both routes
receive 100 tasks. Route A returns 90 answers with 81 correct; route B
returns 60 with 57 correct. B's correctness among returned answers is
95 percent, above A's 90 percent, but its correct-answer fraction among
all assigned tasks is 57 percent versus 81 percent. Neither summary
alone settles utility when correct refusals and harmful wrong answers
have different values. If every unanswered task incurs a stipulated
loss of three and every wrong returned answer a loss of ten, A incurs
120 total loss units and B 150. At loss one per unanswered task, the
totals become 100 and 70. The preferred route changes with the actual
decision consequences. These exact counts show why missingness and
refusal must be classified before ranking an architecture.

A properly specified causal task sometimes requires a refusal or a set
rather than a scalar estimate. Such an output belongs in the correct
answer category if the rubric calls for it. Backend refusal for an
unrelated service-policy reason, inability to parse, and unavailable
scientific method are different categories. They should remain visible
even if a decision maker later aggregates their losses. The prospective
pilot should first establish that the rubric recognizes these outcomes,
then compare useful assistance under that rubric. A subsequent bank
trial can ask whether the resulting architecture preserves meaning and
improves decisions in its own population; the public/generated study
cannot silently stand in for that external question.
\citep{exec14_l6_adapter_exact}

\label{exec14:l6-adapter-end}


% SOURCE docs/release/2026-10-29/source/docs/writing-repair/2026-10-29/src/259a8048adfb1ffa.tex CHAR 0
\section{The implemented boundary}

The historical prototype reviewed in this section implements a local handoff
from a structured business proposal to fixed causal operations. Its preserved
trials and replay define the claims below; the bounded fitted research in
Chapters~\ref{gap29:servicing} and~\ref{gap29:policy-uncertainty} is separate
evidence, not an extension of this translator's tested contract. It accepts an authored fixture or a recorded model
output in a provider-neutral envelope. It does not contact a model endpoint.
The envelope contains the question and its alternative targets, candidate
diagrams, proposed operation names and a provenance record. That record names
the provider, model, revision and generation parameters, retains the exact
prompt, and checks the prompt's SHA-256 identity. The complete proposal receives
its own content identity. A recorded label is provenance supplied by the caller;
it is not independent proof that a particular model generated the bytes.

This boundary is deliberately useful before a provider is selected. It permits
the same semantic checks to be applied to different proposed backends without
making a framework's response format the definition of a causal question.
It also separates two evaluations. One asks whether the controller preserves a
given declaration and result. The other asks whether a model produces good
declarations from unfamiliar business questions. The executed examples answer
the first. They supply no empirical answer to the second.

A proposal can request graph review, graphical identification, discovery or a
moment-based effect calculation. These names do not imply that every operation
is implemented for every input. The fixed service implements syntax review,
a supplied-DAG back-door attempt, and a declared Gaussian moment-inversion
calculation. Discovery remains unsupported in this tested service revision. An unsupported
operation returns that status; it does not invent an empty graph or an effect
of zero. There is no operation for sending a customer communication, altering
an account or approving a bank action.

\section{The three identities of a proposed analysis}

Three identities control the handoff. The proposal identity describes the
complete proposed content and model record. The declaration identity covers the
problem and diagram together, including every target alternative, evidence
reference, assumption, clock and graph revision. Finally, the execution identity
covers prepared requests, the caller's grant, source code, schemas and numerical
environment. Changing a population while keeping the same diagram changes the
declaration. Changing an implementation while keeping the same question changes
the execution. Neither change can silently reuse an old result as if nothing
had changed.

The declaration validator checks structure and selected semantic constraints.
These include whether referenced evidence exists in the declaration, whether
graph endpoints refer to declared variables, whether a DAG has a directed cycle,
whether a proposed history includes a future decision, and whether a policy
uses information before its declared availability. It also flags unresolved
targets, contradicted assumptions, post-treatment adjustment for a proposed
total effect and conditioning a population on treatment descendants. It does
not prove that an asserted DAG is true or fully recognize every MAG, PAG or
CPDAG. A successful declaration review means that a stated problem can proceed
to an appropriately scoped formal analysis.

Evidence identities receive an additional byte check. Source records reference
bounded local files under the documentation tree; symlink escapes and paths
outside that boundary are refused. The actual bytes must match the declared
hash. This catches a changed source or stale citation record. It does not check
whether a source's scientific assertion is true. The full source record,
declaration and proposal are retained as immutable artifacts so that a result
can be traced back to both the claim and the bytes that supported its review.

The external grant is a trusted local caller object bound to the declaration
identity. It lists permitted operations, names the reviewer and supplies the
sampling declaration and any numerical inputs. Putting a reviewer's name into
the model response does not create this object. Requesting a new operation does
not extend its permissions. The grant is an application boundary, not a digital
signature or a bank identity service. Institutional authentication and approval
stores remain requirements for a later deployment.

\section{A narrow projection with explicit refusals}

The translator in these preserved trials supports an unconditional population ATE for one
static binary action, with intervention one and comparator zero. Its horizon
starts at the declared origin. It retains all variables and requires the graph
to account for them. Outcome and target-loss units must agree. It does not
silently convert a dynamic regime, an overlap target, a selected population or
a utility transformation into this static ATE.

For a supported target, the translator builds the fixed operation's scientific
context and places the complete original declaration in its source evidence.
The smaller operation schema carries the treatment, outcome, population, units,
horizon, assumptions and graph. The retained declaration carries the detailed
measurement lineage and numeric decision/availability clocks. Thus a consumer
can reconstruct the original problem even where the narrow operation schema
has no field for every detail. No synthetic or analyst assertion is upgraded to
externally verified truth during this projection.

Candidate graphs are processed individually. For the synthetic offer question,
the input includes an undirected CPDAG and its two directed alternatives. The
controller preserves all three candidate identities. Syntax review returns the
CPDAG's tail endpoints intact. A supported directed graph can receive a
conditional back-door result. Another candidate may be outside the implemented
identification method. The overall record stays ambiguous, and every candidate's
response remains visible. It does not average graph outputs or select the most
convenient result. These alternatives describe possible worlds; the prototype
does not assert their simultaneous truth or estimate probabilities for them.

The distinction between nonidentification and an unsuccessful implementation
attempt is essential here. Failure of the supplied back-door procedure does
not prove that no other valid identification argument exists. Likewise, refusal
to translate the servicing policy does not prove that its value is unknowable.
It records that this controller has no qualified mapping from that richer
target into its fixed service. The full policy, clocks and assumptions remain
available for a future longitudinal adapter.

\section{Representative executed paths}

Six synthetic local paths were executed and exported with their artifact
descriptors and checkpoint records. The ambiguous funding question is blocked
before invoking a tool because its economic target has not been resolved. The
observation example is also blocked: ownership and selection questions remain
substantive. The servicing declaration passes its own chronology checks but is
unsupported by the narrow static translator. These outcomes are useful behavior;
returning a number would answer an unstated substitute question.

The discovery example reaches actual graph-review and identification tools for
each candidate and preserves ambiguity. A separate directed fixture exercises
ordinary conditional execution, interruption and exact replay. It is an explicit
fixture choice of a supplied graph, not a discovered resolution of the original
ambiguous problem. Finally, a weak moment example reaches the effect service
and retains a two-component unbounded set. The mean vector supplied is $(1,0)$,
the covariance matrix is diagonal with entries $0.01$ and $1$, and the declared
sample size is $100$. The resulting set has two unbounded tails separated by a
finite excluded region. The controller preserves both tails and the service's
joint-confidence label together with its Gaussian-approximation limitation.
It neither replaces the union by a finite interval nor interprets the zero
denominator estimate as a finite causal effect.

The numerical example checks propagation of a declared mathematical result.
It does not estimate moments from bank observations or prove that the supplied
covariance is valid under dependence. The original scientific gaps concerning
persistent maps, latent temporal graphs and dependent learners therefore remain
open. A semantic adapter cannot repair a failed inferential law by attaching a
more complete provenance record.

\section{Recovery and adverse inputs}

The local journal stores each completed tool response under a hash of the full
execution identity. Publication is atomic. On resumption, the controller
checks that saved results form a prefix of the planned operations, that their
content hash is unchanged, that each response belongs to the correct request,
and that referenced artifacts still match their bytes. Source verification also
runs before a completed replay. A changed source is therefore rejected even
when a result already exists in the journal.

In the interruption fixture, the controller stops after saving the graph-review
response and before the identification step. Resumption keeps the saved review
and executes only the remaining identification operation. A later complete
replay returns the same record without reinvoking tools. This is a tested
single-host property of the local fixed operations and journal. It does not
establish distributed exactly-once execution, authenticated tamper resistance
against the local operator, or recovery for arbitrary external side effects.
Those stronger claims would require different implementations and trials.

The adverse tests exercise both content and authority boundaries. They reject
unknown operations, injected command or permission fields, nonfinite model
parameters, changed prompt bytes and a live-mode declaration. A retrieved text
fixture tells the system to ignore checks and execute a shell command; the
controller keeps it as source data and performs only the granted graph review.
This verifies the deterministic command boundary. It is not a measurement of a
live model's susceptibility to scientifically misleading text.

Other tests change source bytes, corrupt a checkpoint, alter a target after
the grant, introduce a contradicted assumption, or move the servicing decision
before $L_1$ is available. Missing numerical evidence returns a failed
calculation without an estimate. The recorded output retains the cause of
failure. Together with the inherited problem and operation contract cases,
80 checks pass. This is local conformance evidence, not an accuracy percentage
for unseen banking questions.

The first integration runs also exposed ordinary engineering limits. Some
frozen operation schemas accepted fewer optional fields than the general context
schema. Retaining detailed clocks in the bound declaration supplied a compatible
mapping without dropping them. A 256 MiB address-space allowance was insufficient
for imported numerical dependencies; the executed configuration uses the
existing runtime's tested 1 GiB allowance. Both failed runs are retained in the
evidence. Memory ceilings constrain execution; they do not imply measured
peak memory or comparative efficiency.

\section{What the next evaluation must establish}

The prototype creates a stable place to evaluate actual assistance. A model can
eventually propose a target and diagram, but its output must pass the same
declaration and permission boundaries. An evaluation must compare those proposals
with an ordinary structured workflow and with independent reference judgments.
It must score target meaning, evidence applicability, unresolved ambiguity,
graph classes and downstream effect conclusions. Faster or more fluent reports
are secondary when the intervention, population or clock has changed.

The present records contain zero live model calls. Model tokens and prices
remain unavailable, rather than being reported as measured zeros. All examples
are synthetic, all graph resolutions are explicit fixture assumptions, and no
customer action is authorized. A real comparison still needs exact backend
versions, approved data custody, original benchmark identities, experts for
business and diagram judgments, and a frozen analysis of every requested case.
The controller supplies traceable inputs and outputs for that work; the missing
scientific and institutional evidence cannot be inferred from its successful
execution.


% SOURCE docs/release/2026-10-29/source/docs/writing-repair/2026-10-29/src/c2fbca2ec0b4a3bd.tex CHAR 176

\par The local implementation and historical conformance counts in this chapter are bound to the preserved source, original trials and subsequent six-case replay. They do not constitute new provider execution. \citep{exec14_l6_local}
\label{exec14:l6-end}


% SOURCE docs/release/2026-10-29/source/docs/writing-repair/2026-10-29/src/9782900c4007e2c8.tex CHAR 0
\section{Evaluate the causal contract across the agent's components}
\label{exec18:framework-evaluation}

An agent framework becomes useful to the bank when its components can be
evaluated against a stable decision contract. Consider again an arrears
servicing offer. A plausible explanation, a completed sequence of tool calls
and a correct causal answer are three different outcomes. The system might
produce a fluent justification while conditioning on a variable recorded
after the offer. It might execute the requested estimator successfully while
answering the effect among treated accounts rather than the effect in the
eligible population. Conversely, a scientifically correct answer might be
that the declared observations do not identify the requested effect. A
framework comparison should preserve these distinctions when it chooses
components, evaluation data and metrics.

\subsection{Architecture is a map of responsibilities, not a certificate}

\citet{exec18_framework_agents_survey} organize autonomous-agent architecture
into profile, memory, planning and action, and distinguish architecture from
capability acquisition. This is a useful account of the survey authors'
taxonomy. It is secondary evidence for the individual systems and experiments
that the survey discusses. Its selected March 2025 revision, including a
benchmark table largely describing 2022--2023 work, cannot establish the
current frontier through October 2026. The taxonomy nevertheless helps
specify where a causal assistant's responsibilities reside.

A profile can declare a servicing analyst's remit; it cannot establish the
bank's authority to change an account's treatment. Memory can preserve a
previous graph and its supporting sources; retrieval of that graph does not
make its assumptions current for another product or calendar period.
Planning can propose an adjustment analysis, a sensitivity analysis or a
request for clarification. The plan needs a mathematical applicability
check before execution. An action component can call a causal engine, but
the result must retain the target, graph and observation definitions used
by that engine. These are proposed mappings from the taxonomy to banking
responsibilities, not findings that the surveyed agents already satisfy
them. The state invariant in Equation~\ref{exec14:l6-state} remains the
application's responsibility.

The survey separates subjective human evaluation from objective metrics,
protocols and benchmarks. Task success, similarity to human behavior and
efficiency answer different questions. Its discussion also recognizes the
cost and population bias of human evaluation and the inability of objective
techniques to measure every capability perfectly. In this proposal, a
business reviewer judges whether the question and explanation reflect the
intended decision. A mathematical reference checks whether a result is
correct under explicit assumptions. Neither judgment should silently
replace the other. An uncertain or disputed reference requires an
adjudication status; another model's agreement is not sufficient causal
ground truth. \citep{exec18_framework_agents_survey}

The same survey identifies interacting-module prompt fragility,
hallucination, knowledge-boundary problems in simulations and repeated-call
efficiency. These motivate concrete evaluation cases. A retrieved passage
can contain an instruction inconsistent with the declared target. An
apparently helpful memory can introduce information unavailable at the
decision time. A simulated analyst can exploit knowledge hidden from the
person being represented, making the simulation's task success irrelevant
to the real information constraint. A correct final answer can require many
failed intermediate calls. These are reasons to inspect complete traces and
information availability, rather than infer banking performance from a
general agent benchmark. The cases are proposed causal evaluations;
their frequencies have not been measured here.
\citep{exec18_framework_agents_survey}

\subsection{What a general evaluation library contributes}

\citet{exec18_framework_promptbench} present PromptBench as a modular library
for models, datasets, prompts, adversarial attacks and evaluation protocols.
Their four-stage pipeline initializes the data and model, specifies prompts,
defines input/output handling and metrics, and runs the evaluation. The
selected article reports twelve tasks and twenty-two datasets. Its appendices
describe extension hooks for datasets, models, methods and metrics. These
are relevant design mechanisms for a causal-contract evaluator: a dataset
could contain versioned business questions and admissible responses, while
a custom metric checks target and assumption preservation.

Those mechanisms do not supply the causal reference. For example, a case
asking for an adjustment set needs a declared graph and an appropriate
graphical criterion. A case asking for an effect interval needs a reference
identified set or another justified numerical target. A question whose
meaning changes with the timing of a balance variable needs a clarification
reference. Matching a stored phrase without checking these objects would
reward superficial agreement. A proposed adapter should therefore send a
typed request to the existing scientific boundary and score its typed
response, while retaining the prompt, reference version and execution
trace. PromptBench's extensibility makes that design plausible; no such
adapter is claimed to have been implemented or qualified by this source
inspection. \citep{exec18_framework_promptbench}

The paper's empirical conclusions also argue against a single preferred
prompt. It reports adversarial vulnerability, prompt-engineering methods
that do not outperform the baseline on every dataset, and remaining
weaknesses on selected dynamic tasks. These are attributed historical
results, not a local reproduction or a ranking of present models. Its own
limitations discuss incomplete scenario coverage, metric nuance and the
quality of datasets and prompts. A banking comparison should consequently
retain the unmodified baseline, prespecify perturbations and report
task-specific failures. It should not select a favorable wording after
examining held-out answers and then call that result general robustness.
The paper's historical comparison with other libraries is likewise an
author comparison at that time, not an independent purchasing table.
\citep{exec18_framework_promptbench}

The proposed reuse disposition is reference use of these evaluation
interfaces, followed by conditional integration if a selected implementation
fits the causal request and result schemas. A pinned code revision, package
and dependency rights, environment, input/output behavior and a matched
baseline would have to be inspected before executable adoption. The
article's CC-BY-4.0 notice concerns the paper; it is not evidence of the
software distribution's license. Porting the full library would require
new evidence that its maintenance cost is justified relative to a small
local evaluator. Rejecting every general evaluator merely because it does
not itself prove identification would also discard potentially useful
infrastructure. The appropriate boundary is to reuse orchestration and
measurement only while retaining application-owned scientific references.

\subsection{Completion, correctness and the denominator}

Evaluation must make unsuccessful execution visible without penalizing a
correct scientific refusal as if it were a crash. For a fixed case $i$ and
wording $j$, let $R_{ij}=1$ denote a completed, interpretable response and
$R_{ij}=0$ an unanswered execution. Among completed responses, let
$K_{ij}=1$ denote correctness under the frozen causal contract and
$K_{ij}=0$ an incorrect response. Set $R_{ij}K_{ij}=0$ for an unanswered
case solely for the following accounting identity. With $N$ planned
requests, define
\begin{equation}
 C=\frac{\sum_{ij}R_{ij}}{N},\qquad
 J=\frac{\sum_{ij}R_{ij}K_{ij}}{N},\qquad
 A=\frac{\sum_{ij}R_{ij}K_{ij}}{\sum_{ij}R_{ij}},
 \qquad J=CA .
 \label{exec18:framework-denominators}
\end{equation}
Here $C$ is completion, $J$ is correct responses per planned request and
$A$ is correctness conditional on completion. The last two expressions
involving $A$ apply only when at least one response is completed; $A$ is
undefined if none is completed. These are finite accounting quantities,
not estimates of calibrated causal uncertainty. A correct
nonidentification certificate, an appropriate ambiguity report or a
required clarification can count as completed and correct. The contract
must declare when each is appropriate; indiscriminate refusal cannot
receive credit on answerable cases. \citep{exec18_framework_reference}

An exact authored example makes the difference material. Stipulate twenty
servicing-question families, each with two wordings having identical
outcomes within that family for each system. There are forty planned
requests. The example specifies the answers directly; it contains no model
calls, framework trials, simulated accuracy or bank observations.
\begin{table}[htbp]
\centering
\small
\begin{tabular}{@{}lrrrrrr@{}}
\toprule
System & Correct & Wrong & Unanswered & $C$ & $J$ & $A$\\
\midrule
A & 22 & 2 & 16 & $3/5$ & $11/20$ & $11/12$\\
B & 30 & 6 & 4 & $9/10$ & $3/4$ & $5/6$\\
\bottomrule
\end{tabular}
\caption{Exact stipulated outcomes on forty planned requests in twenty
paired families. Conditional correctness favors A; correct responses per
planned request favor B. No sampling or model-performance claim is made.}
\label{exec18:framework-example}
\end{table}
Thus reporting $A$ alone favors the system that completes fewer cases,
whereas reporting $J$ alone conceals its lower error count. Neither scalar
settles the banking decision. The full counts distinguish reduced coverage
from increased wrong answers, and task strata can reveal whether failures
concentrate on nonidentification, timing or diagram construction.
\citep{exec18_framework_reference}

The paired design also limits uncertainty claims. Two wordings of one
contract are not two new scientific targets. In this example their outcomes
are stipulated to be identical, so treating all forty requests as
independent Bernoulli observations would invent replication. For a future
sampled evaluation, comparisons should retain pairing across systems and
the dependence among wordings of a contract. The sampling law over
contracts, shared model behavior and adjudicator uncertainty determine
whether a cluster analysis, a paired randomization argument or another
method is justified. This fixed example supplies no such sampling law and
therefore supplies no confidence interval.

\subsection{Why the cost of an error can reverse the choice}

Suppose, strictly for illustration, that a correct response has value
$v=10$ USD, a wrong response has loss $p$ USD, and unanswered requests have
zero value before their execution costs. Let $n_{+}$ and $n_{-}$ be the
correct and wrong counts, and let $c$ be total execution and review cost in
USD. Net value per planned request is then
\begin{equation}
 V(p)=\frac{v n_{+}-p n_{-}-c}{N}.
 \label{exec18:framework-value}
\end{equation}
Stipulate total costs $c_A=12$ USD and $c_B=24$ USD over the forty requests.
For $p=20$ USD, $V_A=21/5=4.20$ USD and
$V_B=39/10=3.90$ USD per planned request. For $p=5$ USD,
$V_A=99/20=4.95$ USD and $V_B=123/20=6.15$ USD.
The same stipulated answers thus favor different systems under different
error losses. Every amount is an assumption or exact arithmetic consequence;
none is an observed bank price, benefit or return on investment.
\citep{exec18_framework_reference}

A real servicing evaluation would also need to value justified refusal,
clarification, delay and downstream human work, and distinguish errors
with different monetary consequences. A spurious recommendation to change
a customer's treatment is not interchangeable with an incorrect citation.
Those values must be elicited for a declared decision and information set.
Framework token or request counters contribute to measured execution cost;
they cannot establish these benefits or losses. A low-cost system can be
unattractive if its harmful errors exceed the acceptable rate, even when
the average in Equation~\ref{exec18:framework-value} is positive.

The practical comparison is therefore a paired, versioned set of causal
contracts with explicit completion states, scientific references, trace
inspection and a separately stated economic rule. General agent taxonomies
help locate responsibilities, and general evaluation libraries can help
run that comparison. They leave causal truth, reference construction,
rights qualification and the justified sampling law to be established.
The proposed architecture remains conditional on that work; these two
originals do not establish exhaustive framework coverage or a qualified
banking agent.


% SOURCE docs/release/2026-10-29/source/docs/writing-repair/2026-10-29/src/512e06ac3dc2919f.tex CHAR 0
\section{Transfer an agent design only with its causal question}
\label{exec18:reasoning-transfer}
An agent can give a convincing account of why a payment was recovered and
still choose an uneconomic intervention. The difficulty begins before model
selection: explaining a realised event, predicting a specified mechanism,
identifying an intervention distribution and choosing an action are different
requests. Two contrasting contributions clarify this distinction. A biomedical
position paper proposes combining language models, structured knowledge,
heterogeneous observations and formal causal tools. A philosophical study
instead supplies explicit neuron diagrams and asks a language model to identify
events and their causes. Both offer ideas for the banking product, but their
evidence supports different parts of its design.~\citep{exec18_reasoning_biomedical,exec18_reasoning_neuron}

\subsection{A domain architecture is a research agenda}
The selected initial version of Bazgir and colleagues' biomedical paper,
dated 22 May 2025, is explicitly a position paper. Its proposed workflow
combines literature retrieval, structured knowledge, heterogeneous observations
and external causal methods. It discusses control over consequential actions,
verification, task-specific evaluation, reproducibility and changing model
versions. These are useful architectural requirements. The paper's proposed
personalised and longitudinal applications are motivations, however; it does
not report a new matched experiment qualifying the complete causal agent or a
new identification argument for individual treatment decisions. Its claims
about constituent systems remain secondary reports that require those systems'
own technical originals.~\citep{exec18_reasoning_biomedical}

For banking, the useful transfer is therefore a division of responsibilities.
A language model can propose that a call-centre note contains a hardship
request, that an account event represents a treatment, or that a retrieved
document suggests a causal edge. Each proposal must retain the supporting
observation and its time. The scientific layer must separately decide whether
the proposed variable is measured consistently, whether the intervention is
well defined, and whether the available information identifies the requested
contrast. A graph of relationships extracted from documents does not by itself
supply a causal graph for the population or distinguish a predecision
covariate from a consequence of an earlier action. Calling an external
estimator can implement an approved analysis; it cannot make those scientific
premises true. This is the proposed banking interpretation of the architecture,
not a claim that the biomedical paper establishes the banking premises.

The source also illustrates why retrieval must preserve citation identity.
Its discussion of multimodal causal discovery attaches reference 7 to a
MATMCD passage, while its bibliography identifies reference 7 as Trialscope
and reference 17 as the multimodal-agent work. This internal mismatch does
not establish that the surrounding scientific statement is false, but it
prevents that reference number from serving as a verified technical path.
The product should preserve the passage and the mismatch, find the intended
original, and verify the claimed method there before treating it as available
support. An agent-written explanation of the mismatch would remain a
resolution proposal until the original is inspected.~\citep{exec18_reasoning_biomedical}

For a proposed servicing study, public or synthetic observations can test
whether the system preserves times, variables, graph restrictions and refusals
through this chain. Such a study can demonstrate the merit of building a
prototype without requiring bank deployment first. It cannot establish the
truth of a particular bank's undocumented assignment or audit process. The
proposal should therefore state what the public experiment would decide,
which scientific assumptions it stipulates, and what later institution-specific
evidence would change the decision. This boundary permits useful prototype
work while keeping the proposed transfer conditional.

\subsection{A supplied causal mechanism changes the evaluation task}
The selected Vervoort--Nikolaev manuscript has an internal dateline of
30 April 2025 and an arXiv submission of 17 June 2025. It presents textual
transcriptions of neuron diagrams, with stipulated activation, inhibition
and time-order rules. Its two questions ask whether the target event occurs
and what causes that occurrence or nonoccurrence. The model receives the
mechanism; it is not asked to identify a causal graph from a sampled
observational distribution. Performance on this task therefore tests a
particular form of reasoning about supplied mechanisms. It cannot be read as
performance on statistical discovery or intervention identification from
unrestricted data.~\citep{exec18_reasoning_neuron}

The proposed labeling rule, DEF-1, begins with a counterfactual comparison
but introduces additional treatment of bifurcating paths. In selected cases it
retains factual blocking on off-path routes even when the candidate cause is
changed. The authors explicitly describe this as not fully implementing the
ordinary removal of that event throughout the mechanism. They also limit
the rule to a class of diagrams, discuss an ambiguous case requiring a
possible extension, and leave its relation to other causal theories open.
For the banking product this is a candidate attribution convention with a
declared domain, not a universal replacement for an intervention operator.
Implementing a library operation under this convention would require its
own formal specification and conformance tests. The example below does not
implement or validate DEF-1.~\citep{exec18_reasoning_neuron}

The evaluation evidence requires equally careful interpretation. The paper
describes the experiments as illustrations, explicitly distinguishes them
from large statistical studies, and reports that the displayed early tables
came from a single run without feedback. Its 25 selected diagrams also
helped develop the proposed definition. The prose reports 13 of 25 fully
correct answers and two partially correct answers for the selected GPT-4
exercise; these are historical, author-reported counts under that rubric,
not an estimate of general causal competence. The later model comparison
changes the prompt to request short answers and initial causes. Its result
cannot isolate a model-change effect from the changed evaluation contract.
No current model ranking follows from these experiments.~\citep{exec18_reasoning_neuron}

Even a small benchmark needs a faithful transcript and an explicit grading
unit. In the first diagram's recorded answer, the explanation assigns the
blocking of the backup path to A although the prompt assigns it to C; the
answer is nevertheless marked correct. The diagram 17 transcription asks
whether G occurs and then asks for the causes of E. These are local
transcript and target inconsistencies, not a basis for recomputing all scores
from extracted text. They show why an exact-output evaluator must distinguish
a correct final event, a correct cause list and a faithful explanation. A
reuse decision should preserve the originals, record disputed labels and
adjudicate the target before admitting the case. Diagram images, underlined
table annotations and full score reproduction need their own inspection.~\citep{exec18_reasoning_neuron}

\subsection{A servicing route with backup: attribution versus value}
\label{exec18:reasoning-example}
Consider a deliberately stipulated servicing mechanism. Let $A\in\{0,1\}$
denote availability of a backup contact, and $C\in\{0,1\}$ activation of a
priority service. At the next stage, $B$ records whether the backup actually
runs and $D$ whether the priority route delivers. Recovery $Y$ follows the
Boolean equations
\begin{equation}
 B=A(1-C),\qquad D=C,\qquad Y=\max\{B,D\}.
 \label{exec18:reasoning-mechanism}
\end{equation}
These are authored assumptions, not estimated account behavior. They encode
perfect success of an executed route and suppression of the backup while
priority is active. In the factual state $A=C=1$, the realised path is
$D=1$, $B=0$, $Y=1$. A faithful narrative can identify the executed priority
route. Under an intervention setting only $C=0$, however, the structural
equation for $B$ is still active: $B=1$ and $Y=1$. Removing priority alone
does not remove recovery in this state.~\citep{exec18_reasoning_reference}

If a different comparison also forces the backup to remain off, the outcome
is instead $Y=0$. With the subscript describing the intervention,
\begin{equation}
 Y_{do(C=0)}\bigm|_{A=1}=1,\qquad
 Y_{do(C=0,B=0)}\bigm|_{A=1}=0.
 \label{exec18:reasoning-interventions}
\end{equation}
The second expression changes two mechanisms. It may be relevant to a
declared attribution question, but it is not the outcome of disabling only
the priority service. A result object must retain which comparison was
requested. Replacing an intervention by a compound action that holds a
downstream route fixed silently changes the scientific question. Conversely,
a zero effect of disabling one route does not erase the fact that the route
was executed. The explanatory and intervention answers can differ without
logical contradiction.

To make the distinction economic, suppose a recovery is worth USD~200,
priority costs USD~40 when activated, and an executed backup costs USD~10.
The stipulated net value per case is
\begin{equation}
 N=200Y-40C-10B\quad\text{USD}.
 \label{exec18:reasoning-money}
\end{equation}
At $A=1$, priority produces USD~160 net and the backup-only route USD~190.
Priority's incremental net value is therefore minus USD~30 although it
provides the factual recovery path. When $A=0$, priority gives USD~160 while
no priority gives zero. If backup availability is not known before action
but its stipulated probability is $p$, averaging these two states gives
\begin{equation}
 \mathbb{E}[N_{do(C=1)}-N_{do(C=0)}]
 =160-190p\quad\text{USD per case}.
 \label{exec18:reasoning-threshold}
\end{equation}
Under these assumptions priority has positive incremental value precisely
when $p<16/19$, with indifference at equality. If availability is actually
observed before action, the appropriate conditional decision differs from
this population decision. The information clock must therefore accompany
the value calculation. Exact enumeration of all four $(A,C)$ states and
rational checks at $p=0,1/2,16/19,1$ reproduce the result. Neither the
mechanism nor $p$ nor the monetary inputs were learned from bank data; these
calculations establish an internal distinction, not bank ROI.~\citep{exec18_reasoning_reference}

\subsection{Admit the question and the rubric together}
A useful public evaluation can now separate tasks that a single accuracy
score would obscure. One task checks propagation of a supplied mechanism.
Another checks the stated attribution convention. A third asks for a
specific intervention and requires the correct action set. A fourth asks
for a value comparison under a stated predecision information set. An
answer that computes the compound intervention correctly but returns it for
the single-action query fails that contract. An answer that requests the
missing attribution convention may be scientifically preferable to an
unqualified list of causes.

The next prototype design should freeze the target, intervention semantics,
grading rubric and explanation requirements before evaluation. New graph
families should be held apart from those used to design the rules; word
changes alone do not make a familiar mechanism a new causal case. A public
corpus can retain disputed examples with an adjudication flag, rather than
silently treating one philosophical judgment as uncontested truth. Matched
comparisons should retain the same target, prompt conditions, tool access
and output budget, while reporting invalid, unsupported and incomplete
outputs separately. This is a proposed route to credible prototype evidence.
It leaves open the broader coverage of the literature, executable package
admission, live performance and the validity of institutional causal premises.


% SOURCE docs/release/2026-10-29/source/docs/writing-repair/2026-10-29/src/228dc542ffea3793.tex CHAR 0

\section{Memory benchmarks: distinguish provenance, correctness and causal meaning}
\label{sec:exec23-recon-review}
A bank analyst may retract a transaction classification after a later fraud
alert, revise a business definition, or discover that two reports share the
same source. An agent should update conclusions that depend on the changed
premise while retaining conclusions supported independently. Public generated
cases can test this behavior before a costly bank prototype. The benchmark
must nevertheless distinguish a correct derivation within its authored world
from truth about a bank and from an identified intervention effect.

Arya's RECON preprint provides a relevant case design: 24 long case files,
1,604 questions, three domains and six tasks involving chains, invalidations,
source conflicts, alternative timelines, temporal constraints and fact
retrieval. A deterministic skeleton establishes provenance dependencies;
LLMs realize its surface narrative under validation constraints. These are
source-described design and performance claims, not executions in this
workspace. \citep{exec23_recon_original}

\subsection{An authored dependency graph is a useful but different object}
The generator's edge records that one entry establishes an attribute required
by another. Reachability and invalidation determine answers in that authored
world. This can be an excellent test of evidence custody and revision without
being a structural causal model for an external bank. In particular, changing
a timestamp and recomputing reachability is not automatically an intervention
on a data-generating mechanism. The operation and its counterfactual semantics
must be specified. \citep{exec23_recon_original}

A familiar finite distinction can be stated without any learner. Let $U$ be
fair binary and consider two models. In the first, $A=U$ and $Y=A$; in the
second, $A=U$ and $Y=U$. Both give observational probability $1/2$ to
$(A,Y)=(0,0)$ and $(1,1)$. Their intervention contrasts are respectively one
and zero. A record graph can preserve the observed statements in either
world. Perfect retrieval of those records therefore does not identify which
intervention contrast is appropriate. This is an authored example, not a
reported RECON defect. \citep{exec23_recon_reference}

The source separates structural checks, semantic audits and human review.
That separation is valuable, but a schema check cannot establish that every
narrated qualification retains its intended meaning. A name, identifier and
timestamp may all survive while a negation or exception changes the statement.
Human disagreement cannot be assigned exclusively to reasoning difficulty
merely because a structural validator passes. A banking adaptation should
retain the exact skeleton, narration, validation decisions, rejected cases
and independently adjudicated semantic discrepancies.
\citep{exec23_recon_original}

\subsection{Scoring changes the apparent comparison}
RECON's main evaluation removes 190 of 1,604 questions after a closed-book
rule involving three solvers, leaving 1,414. That defines a selected test
population. A correct closed-book answer can arise from guessing or task
structure; the filter is not direct proof of training contamination. Selection
also uses particular models and abstention behavior. Results should preserve
both population identities and the selection flags before generalization.
\citep{exec23_recon_original,exec23_recon_reference}

For six-option single-answer questions, correct, wrong and abstained answers
score $1,-1/5,0$. Uniform guessing has expected score
\[
 \frac16-\frac56\frac15=0,
\]
but strict accuracy $1/6$. Abstention has score and accuracy zero. An agent
correct on one tenth of questions and wrong on the rest has accuracy $1/10$
and expected score $-2/25$, while an always-abstaining agent has zero for
both. The accuracy ranking and score ranking reverse. This is a designed
incentive distinction, not a calculation error. \citep{exec23_recon_reference}

The benchmark additionally uses multi-option weighted scores, exact integer
matches, ordering correlation and the average of two binary LLM judgments for
free-form responses. A macro score therefore depends on the fixed task and
format mixture, and two cross-family judges do not create independent truth
merely by agreeing. A product comparison should retain the underlying answers,
format scores, strict correctness, abstentions, parser failures and judge
votes. It should price these outcomes using a declared business loss rather
than silently interpreting benchmark points as dollars.
\citep{exec23_recon_original}

\subsection{Reported human denominators need reconciliation}
\label{eq:exec23-recon-denominators}
The original describes 200 human-validation questions with three annotations
per question and reports majority agreement with gold of 86.3\%. For an
unweighted binary majority result over 200 items, percentages move in
increments of $100/200=0.5$ points. No integer success count rounds to
86.3\% at one decimal. A rate of $518/600$ does round to that number, but
uses annotation judgments rather than the stated question-majority
 denominator. It is a possible explanation to investigate, not a recovered
fact. Similar questions apply to several reported quality percentages.
\citep{exec23_recon_original,exec23_recon_reference}

The separate human baseline states 60 questions, ten per task, with final
answers determined by majority vote. An unweighted binary total over 60
cannot round to 63.0\% at one decimal; ten-item task rates cannot be 78\%,
58\%, 45\% or 37\% under strict majority correctness. These checks do not
establish fabricated results. Unreported weights, exclusions, averaging or
annotation-level summaries could explain some values, but they would change
the stated estimand and denominator. Exact item/annotator records, aggregation
code and missingness rules are needed before using the human comparison as
quantitative investment evidence. \citep{exec23_recon_original,exec23_recon_reference}

Human annotators receive question-specific evidence packets, whereas other
systems receive the full narrative, retrieved excerpts or structured records.
The paper acknowledges this asymmetry. The human comparison cannot isolate
reasoning ability from evidence selection and representation. Likewise, the
source's Oracle is an answering LLM given structured truth, not an algorithm
that directly returns the generated answer. Its weaker temporal-fact result
shows why it is not a universal mathematical upper bound on every system.
\citep{exec23_recon_original}

\subsection{Retrieval, dependence and full cost}
The reported support-recall diagnostic searches for verbatim composite entry
headers. The source records imperfect header uniqueness and excludes two
systems whose ingestion paraphrases away those headers. Header presence does
not establish that the full fact, qualification or independent source reached
the answerer; header absence can coexist with a faithful paraphrase. The
observed hit/miss accuracy contrast also conditions on a potentially easier
subset. A paired intervention supplying or withholding matched evidence would
be needed to attribute an accuracy change to retrieval itself.
\citep{exec23_recon_original}

The source reports a question-level paired bootstrap. Its appropriateness
depends on whether the target is the fixed question collection or new
independent case files. An authored counterexample has two independent case
signs, each repeated in ten question contributions. The mean over all twenty
questions equals $(Z_1+Z_2)/2$ and has variance $1/2$. Treating all questions
as independent gives $1/20$, ten times smaller. This does not refute every
within-benchmark summary; it demonstrates why a new-case target needs the
case dependence, task mixture and effective independent units stated.
\citep{exec23_recon_reference}

The reported narrator-family tests have small case counts and include a
13.3-point difference with a nonsignificant test. Nonrejection does not
establish equivalence or rule out a practically important effect. A prospective
adaptation should preregister an equivalence margin or report uncertainty
without promoting a nonsignificant result into absence of bias.
\citep{exec23_recon_original}

Finally, prompt/completion tokens per answer are not total workflow cost.
Ingestion, embedding, reranking, query rewriting, indexing, storage, retries,
judges and failed runs must be included. The paper specifies service wait
intervals of 30, 60 or 600 seconds for different routes; elapsed waiting alone
is not an observed completion or durable-readiness receipt. The proposed
CLI/MCP worker should bind input and index versions, completion state,
retrieval artifacts and cancellation/restart evidence before declaring a
comparison valid. No service has been exercised by this source review.
\citep{exec23_recon_original}

RECON is retained as a public/generated test-design reference for revision,
provenance and compositional reasoning. Its original code, dataset versions,
filtered row identities, human denominators and complete costs remain
acquisition and qualification obligations. A future implementation can compare
matched evidence and answerers under a declared loss; it cannot convert
benchmark provenance into bank truth or causal identification by naming the
graph causal. Native performance, bank benefit and general coverage remain open.


% SOURCE docs/release/2026-10-29/source/docs/writing-repair/2026-10-29/src/d6ab2a32bb923e37.tex CHAR 0
\section{Specialist state, typed interfaces and learned latent collaboration}
\label{sec:exec24-collaboration}
A servicing analyst may ask an assistant to compare repayment policies while a specialist holds a transition law, uncertainty and diagnostic failures. Returning only a preferred action can destroy information needed to explain a refusal or change the loss function. Returning every hidden activation does not establish an identified intervention effect. The question is which information the decision requires and how its interface preserves it. LLAMIA offers a concrete hypothesis about specialist--language-model communication. We inspect the latest version listed before our cutoff, the 46-page v2 of 2 September 2026, while preserving v1. We retain the architecture as a prospective experiment, with unresolved reporting contracts distinguished from our own exact examples.~\citep{exec24_collaboration_original,exec24_collaboration_reference}

\subsection{Information, representation and a substantive alternative}
The paper couples a frozen chess specialist with a language model. A three-layer projector maps specialist activations to 32 continuous tokens in the language model's embedding space. Both principal variants receive verbalized engine outputs; the latent variant additionally receives projected activations. The first training stage aligns the projector while freezing the language model; the second jointly trains both with task rewards. Human-move prediction, puzzle difficulty and interest, move explanation and commentary form the chess suite; Go move prediction supplies a second setting. These public-data designs expose a specified environment before expensive bank work. They do not identify intervention effects on repayment or funding.~\citep{exec24_collaboration_original}

The comparison changes information, representation and a learned interface together. It does not compare every text interface with every latent interface. Top-move summaries, complete typed policy/value records and task-specific statistics are different alternatives. The paper reports probe, template, raw-encoder and shuffled-state controls. Weaknesses of those implementations cannot prove that sufficient information is inherently inexpressible as text or that hidden activations suffice for a new bank decision. The paper acknowledges architecture, specialist-strength and projector-compatibility confounding in one ablation.~\citep{exec24_collaboration_original}

Let $H$ be specialist state, $a\in\mathcal A$ an action, $L(a,Y)$ a loss, and $B=b(H)$ the interface object. For finite action sets and integrable losses,
\begin{equation}
 R_H=E\!\left[\min_{a\in\mathcal A}E\{L(a,Y)\mid H\}\right],\qquad
 R_B=E\!\left[\min_{a\in\mathcal A}E\{L(a,Y)\mid B\}\right].
 \label{eq:exec24-collaboration-sufficiency}
\end{equation}
Every $B$-measurable policy is $H$-measurable, so $R_H\leq R_B$. Equality is available if an optimal action based on $H$ is measurable with respect to $B$; a stronger sufficient condition makes every conditional action loss a function of $B$. This is an information statement under a stipulated law and an ideal decision rule, not a learning or causal-validity guarantee. A finite-precision record can be serialized injectively into text. Existence of that serialization does not establish context efficiency or learnability. A learned projection can also lose relevant distinctions.~\citep{exec24_collaboration_reference}

In an authored example, equally likely states $H=-1,+1$ have default risks $(0.02,0.04)$ and $(0.04,0.02)$ under actions $(0,1)$. Loss given default is USD~10,000; the intervention has no additional cost in this narrow stipulated law. Its advantage over action zero is USD~$-200$ in one state and USD~$200$ in the other. Full information permits USD~100 expected savings relative to always choosing zero. A constant summary permits none; a typed one-bit state identifier permits the same USD~100 as full information. The exact reference enumerates this law. It demonstrates one lossy summary and a sufficient textual alternative, without estimating bank benefit or proving that a hidden vector is necessary.~\citep{exec24_collaboration_reference}

\subsection{Reconcile results before using a ranking}
The v2 tables disagree. Table~14 gives the 14B latent Wild scores $57,54,60$, averaging $57$, while headline Table~18 gives $49$. Table~15 gives rationale components $44.0,56.6$, averaging $50.3$, while Table~18 gives $45.8$. Table~17 gives difficulty correlation $0.67$, while Table~18 gives $0.71$. Layout and raw extraction agree on these entries. Different checkpoints or subsets might explain them; the inspected original does not supply that mapping. We do not select a preferred value or silently repair the source. Broad expert-superiority wording also needs qualification against reported MAIA scores $53$ for LLAMIA and $55$ for the expert.~\citep{exec24_collaboration_original,exec24_collaboration_reference}

The data account specifies 100 held-out commentary games and 5,000 puzzles; result captions instead describe approximately 300 games and 1,000 puzzles. One solved-puzzle entry remains a placeholder. Exact row identities, runs, exclusions and aggregation rules are needed for each score and interval. These discrepancies require reconciliation, not an allegation of fabrication. Ratios of correlations, BLEU scores and judge ratings are not ratios of monetary benefit.~\citep{exec24_collaboration_original}

The human study illustrates a coherent denominator distinction. Twelve raters evaluate 15 positions, producing 180 crossed judgments. A count of 130 favorable responses yields $130/180\simeq72.2\%$ overall and $130/162\simeq80.2\%$ after 18 ties are removed. Both can be correct. A pooled non-tie rate need not equal the average of per-rater rates when ties differ. Neither denominator creates 180 independent raters. For an authored model $Z_{ij}=U_i+V_j$ with independent mean-zero unit-variance rater and position effects,
\begin{equation}
 \operatorname{Var}(\overline Z)=\frac1{12}+\frac1{15}=\frac3{20},\qquad
 \operatorname{Var}_{\rm iid}(\overline Z)=\frac2{180}=\frac1{90}.
 \label{eq:exec24-collaboration-crossed}
\end{equation}
The ratio is $13.5$. This does not recalculate the paper's uncertainty; it shows why new-rater, new-position and fixed-panel targets need different arguments. Reusing a judge for reward and evaluation also calls for independent adjudication and reward-overfitting diagnostics.~\citep{exec24_collaboration_original,exec24_collaboration_reference}

\subsection{State refresh, training and cost contracts}
Section~J.3 refreshes state tokens after board mutations and excludes read-only calls, including \texttt{get\_policy}, from re-encoding. Section~J.4.3 says \texttt{get\_policy} performs a fresh forward pass and injects 32 tokens. A consumer cannot infer which rule produced a run. An integration must bind state and query versions, specialist checkpoint, activation layer and shape, projector revision, refresh trigger and stale-result policy. Parallel reads need a common snapshot or tagged versions; sequential writes alone do not bind delayed reads to the correct state.~\citep{exec24_collaboration_original}

The narrative claims no human-move supervision, while the later training account includes 500,000 positions paired with human moves and a rank reward. No direct supervised fine-tuning and no human-move information are different claims. Alignment, labels and reinforcement rewards need separate provenance. Joint projector training also requires a precise old-policy likelihood: stored rollout features and features recomputed with a changed projector are not automatically the same conditioning history. Formulas alone do not settle this implementation. This is a source/code requirement, not an observed failure.~\citep{exec24_collaboration_original}

Cost neutrality is not established by returned-token counts. Table~2 gives 8B training costs of 16.8 versus 15.0 GPU-hours: 12\% overhead, inconsistent with approximately 6\% at every scale. At 14B, reported calls and return lengths yield $2.9(150)=435$ text tokens versus $1.9(182)=345.8$ latent-plus-text return tokens. These products reproduce narrow return accounting only. They exclude other prompts, output, repeated attention to history, specialist computation, projection and training. Continuous embeddings are not automatically executed or billed like vocabulary tokens. Five million alignment pairs, the multi-GPU configuration and step counts require elapsed and aggregate GPU-time with explicit inclusion rules.~\citep{exec24_collaboration_original,exec24_collaboration_reference}

\subsection{A bounded public comparison and integration decision}
Freeze public or generated tasks, splits, response laws, specialist/model/projector versions, loss functions and budgets. Compare top-$k$ prose, complete typed records, a sufficient statistic where available, latent-only access and typed-plus-latent access. Preserve no-specialist, shuffled/missing-state and stale-result controls; include failures and refusals in the denominator. Match training resources when attributing interface effects and report cost differences when matching is impossible. A game demonstration can establish feasibility; a generated servicing environment tests banking semantics under declared assumptions. Neither supplies institutional validation.~\citep{exec24_collaboration_original,exec24_collaboration_reference}

For CLI/MCP design, the primary deliverable remains a typed auditable result containing target, units, assumptions, uncertainty and refusal status. A latent extension should reference immutable artifacts and state identities; a custom model host additionally needs tested embedding injection and specialist-activation access. A text tool response alone does not demonstrate that path. The original suppresses refusal and mandates numerical guesses in task prompts. We reject that rule for bank inference: unidentified effects, unsupported regimes and stale state require typed refusals. Model, corpus, code and dependency rights also precede adoption. Retain latent communication as a conditional experiment and complete typed records as a substantive comparator; defer native integration until reporting and execution contracts are reconciled. Native performance, bank benefit and general coverage remain open.~\citep{exec24_collaboration_original,exec24_collaboration_reference}


% SOURCE docs/release/2026-10-29/source/docs/writing-repair/2026-10-29/src/435d2a851eef8e58.tex CHAR 0
\section{Agent forgetting as an observable contract}
\label{sec:exec24-unlearning}

A bank that removes an obsolete customer record from an agent's retrieval store
still has to ask what remains in a cached tool reply, an analyst summary, a prompt,
and a fitted model. These are different information surfaces. A correct answer to
an intervention query does not establish authorized information custody. K-Bench
supplies a useful public experimental design for this distinction. Yu and
colleagues observe tool-using agents across parametric, context, textual retrieval
and structured retrieval configurations, measuring disclosure alongside retained
functionality and protocol degeneration. We use the September 14, 2026 v2 original
and the public repository revision available before the October 7 cutoff. Its
findings motivate a prospective evaluation contract; they do not establish that
our causal library has acquired a working unlearning component.
\citep{exec24_unlearning_original,exec24_unlearning_software}

\paragraph{What the observer measures.}
Let $q$ denote a query about a specified entity and field, $v_q$ its target value,
and $T_c(q)$ the text emitted on observed channel $c$. The paper considers six
channels: emitted reasoning or scaffold text, tool arguments, the wider tool
observation, retrieved identifiers, an answer and a separately elicited summary.
These observations do not presume access to an inaccessible internal reasoning
state. The summary is a new elicitation, rather than simply another view of the
same ReAct trace. Define the declared matcher $d$ and observed-channel set
$\mathcal C_q$ explicitly:
\begin{equation}
 L(q)=\max_{c\in\mathcal C_q}d(v_q,T_c(q)),\qquad
 \widehat R=N^{-1}\sum_{q=1}^{N}L(q).
 \label{eq:exec24-unlearning-observer}
\end{equation}
The fixed observer detects a target present in any single channel. It is not an
adaptive sequence of attacks, a proof of removal from weights, or a universal
information-flow guarantee. Its public matcher canonicalizes text and uses
substring containment, with field-specific date and address variants. Opaque
retrieved identifiers are not automatically counted as disclosures merely
because retrieval occurred. The graded metric additionally credits recovered
target tokens and takes the largest channel severity for each query.
\citep{exec24_unlearning_original,exec24_unlearning_software}

Channel denominators require care. The public scorer omits failed summaries
from that channel's mean and omits an answer only when there is no observed
answer evidence. The aggregate still averages over queries. Thus a conditional
answer-channel rate need not be bounded above by an aggregate with a larger
denominator. In an authored four-query example, two queries carry an answer and
one answer discloses the target. The conditional rate is $1/2$, while the
all-query rate is $1/4$ if no other query discloses anything. Apparent numerical
inconsistency must first be examined against these populations. The paper's
discussion of omitted halted answers and the current raw-answer fallback also
make the scorer revision part of the endpoint.
\citep{exec24_unlearning_software,exec24_unlearning_reference}

\paragraph{A union of detectors is not joint reconstruction.}
The original explicitly excludes an attacker who assembles partial values from
different channels. A minimal authored text example has target \emph{alpha beta},
one channel containing \emph{alpha} and another containing \emph{beta}.
Complete-value matching returns zero on both channels; the largest token recall
is $1/2$; concatenation recovers the entire stipulated target. The stronger
distinction is information-theoretic. Let $S$ and $U$ be independent fair bits,
with $Z_1=U$ and $Z_2=S\mathbin{\oplus}U$. Enumeration of the four equally likely
states gives
\begin{equation}
 H(S\mid Z_1)=H(S\mid Z_2)=1,\qquad H(S\mid Z_1,Z_2)=0.
 \label{eq:exec24-unlearning-joint}
\end{equation}
Each isolated channel leaves a full bit of uncertainty; the pair determines
$S$. These are project derivations, not a reproduction of the benchmark's
matcher or a theorem attributed to its authors. If a stipulated bank decision
must choose the bit correctly and an error costs USD 200, either isolated
channel leaves expected loss USD 100, while their combination permits zero.
This is a designed value-of-information illustration, not a measured privacy
cost, price or banking result. An approved reconstruction observer and its
allowed side information must accompany a disclosure claim.
\citep{exec24_unlearning_original,exec24_unlearning_reference}

Entity coverage is a second dimension. A detector keyed to a request for entity
A's value can report no leak when the reply instead discloses entity B's value.
If B also belongs to the forget set, that is a missed forget-set disclosure even
though the A-specific endpoint is computed correctly. The paper acknowledges
this limitation and reports that counting other forget entities changes the
retrieval results for its input-corruption reference. A bank evaluation should
report queried-entity, any-authorized-forget-entity, third-party and
joint-reconstruction endpoints separately. The broader endpoints require a
declared ground-truth corpus and access policy. Synthetic linked entities and
deliberately split values suffice to test the initial contract.
\citep{exec24_unlearning_original,exec24_unlearning_reference}

\paragraph{Reduction, utility and collapse are separate decisions.}
The binary reduction label, graded scalar and eligibility gates answer different
questions. The paper's scalar has the form
\begin{equation}
 K=(1-\overline R_f)(1-|\Delta_r|)_+
   \{1-\max(0,D_f-D_{f,0})\}_+,
 \label{eq:exec24-unlearning-score}
\end{equation}
where $\overline R_f$ is graded forget-set severity, $\Delta_r$ the graded
retain-set change, and $D_f$ and $D_{f,0}$ candidate and baseline degeneration.
The public CLI separately requires a retain ratio of at least $0.80$, added
degeneration at most $0.20$, and no candidate split reaching $0.50$ absolute
degeneration; the untreated baseline is exempt from the terminal label.
A positive scalar cannot itself establish eligibility. Under the authored
values $D_{f,0}=1/2$, $D_f=1$, zero forget severity and zero retain shift,
the scalar is $1/2$ although the terminal gate fails. This makes the baseline
dependence explicit without claiming that an empirical cell has these values.
\citep{exec24_unlearning_original,exec24_unlearning_software,exec24_unlearning_reference}

Two definition discrepancies remain visible. Equation (3) prints a signed
retain difference below a positive tolerance, whereas the next paragraph
requires its absolute value below that tolerance. A change from $4/5$ to $1/5$
gives $-3/5<1/20$ but $|-3/5|>1/20$; only the latter captures the stated retention
intent. Separately, page 13 calls a significant reduction below twofold a
measured failure. The inspected public binary classifier returns a reduction
label for such a case. For stipulated baseline $3/5$, candidate $2/5$ and
adjusted $p=1/100$, its visible branch yields a $1.50$-fold reduction label,
while the printed twofold condition fails. This is a static-contract finding
with an authored input, not execution of downloaded code or proof that a
published empirical row traversed that path. The CLI calls that classifier
and also reports eligibility separately. An adopted evaluator needs one
named rule set and explicit migration of old reports.
\citep{exec24_unlearning_original,exec24_unlearning_software,exec24_unlearning_reference}

Protocol health differs from correct service. The current scorer inspects raw
text and the first answer payload rather than trusting a terminal label alone.
An unhealthy trajectory can still disclose information, so its answer evidence
remains relevant. A bare direct reply is recovered as answer evidence only
under specific no-tool/no-thought conditions. Valid answer formatting does not
prove correctness, causal identification, appropriate refusal or customer
benefit. The bank contract needs those outcomes alongside health and disclosure,
retaining failure reasons and permitted trace provenance.
\citep{exec24_unlearning_software}

\paragraph{Experimental units and causal interpretation.}
The observer tests a declared agent configuration. It does not identify a
customer-intervention effect or establish that storage location alone caused
a cross-configuration difference. The parametric lane uses an injected adapter
or merged target, while nonparametric lanes can use the unmodified base.
Protocol amendments explain this choice through failures of the original
adapter-loaded design. That explanation is valuable feasibility evidence,
but model state and storage location then vary together. A prospective bank
study should either cross those factors with viable controls or describe the
contrast as a bundled configuration change. The oracle-gated input-corruption
reference also replaces learned components and retains the secret behind its
gate; it is not evidence of verified weight erasure.
\citep{exec24_unlearning_original,exec24_unlearning_software}

The paper pools paired query outcomes across seeds for exact McNemar tests and
separates forget and retain multiple-testing families. The code samples queries
from entity-field records, pairs baseline and candidate within seed, then adds
discordant counts. The public CLI checks uniqueness and equality of query
identifiers within each compared seed and restricts the baseline to the
candidate's seeds. These are meaningful alignment controls. They do not establish
independent discordant signs when entities, prompts, training data or runs are
shared. The standalone verdict routine pairs intersections of identifiers and
deduplicates query IDs by keeping the last row; that weaker path must not be
silently substituted for the CLI's cohort checks. Its topology aggregation
averages bootstrap interval endpoints across seeds, an explicitly approximate
construction whose coverage is not proved by the existence of the code.
\citep{exec24_unlearning_original,exec24_unlearning_software}

For an exact authored dependence witness, take 200 independent fair Bernoulli
entity outcomes and repeat each three times. The mean of 600 rows equals the
mean of 200 independent values, so its variance is $1/800$, not $1/2400$ from
treating all rows as independent. This is not an estimate of K-Bench's actual
variance. It shows what the experiment manifest must rule out or model. A public
prototype should preserve entity, query, seed and run IDs, distinguish
finite-corpus from new-entity targets, and predeclare appropriate uncertainty.
Splitting multiplicity families does not repair an invalid input test; changing
gates after observing results needs a dated amendment and a new evaluation
population.~\citep{exec24_unlearning_reference,exec24_unlearning_software}

\paragraph{Reuse, provenance and cost.}
The repository is a useful candidate reference for channel records, health
diagnostics, cohort validation and offline report exchange. Its direct code
license is MIT, whereas the synthetic corpus declares CC BY 4.0 and model and
retrieval assets carry separate terms. Main requirements pin versions but are
not the complete original lock; the compute guide explicitly leaves an
OpenUnlearning PEFT version unresolved. The reproduction guide describes
approximately 500 GPU-hours for the whole study and about 18 GPU-hours for
selected released targets after preparation, identifying paper-only figures
those targets cannot rebuild. These author estimates cannot become a bank
budget without training, acquisition, storage, inference, failure, review and
accepted-output accounting. A conditional original-Python worker is a prospective
evaluation option, with a small typed transcript boundary for our CLI/MCP.
A whole-library port has no demonstrated advantage. Automatic adoption is
withheld pending rule reconciliation, environments, rights review and registered
conformance.~\citep{exec24_unlearning_software}

Artifact custody is part of that boundary. The inspected local-model fingerprint
hashes configuration bytes and weight filenames and sizes, but not weight
contents. Two different four-byte payloads have the same name and size yet
different content hashes, so that fingerprint cannot establish byte identity.
The source separately resolves a hosted repository revision, a stronger identity
than its README's shorthand main-revision description; API model IDs are another
kind of identifier. Our receipt should bind actual weights or hosted immutable
revisions, tokenizer, prompt, dataset, scorer, dependencies, observation policy
and experiment declaration. A document hash or successful exit cannot stand in
for those individual identities.~\citep{exec24_unlearning_software,exec24_unlearning_reference}

The ecological-data description remains unresolved. Pages 41--42 describe real
names and birth dates from a converted external dataset, while page 49 says all
personally identifiable information is synthetic. The pinned conversion script
calls its input synthetic and the dataset sheet says the external entities are
not redistributed. These support a provenance question, not a misconduct finding
or a license to choose the more convenient account. The exact external data
version, conversion output and source rights are needed before making a
real-data validation claim. Initial bank-facing experiments can use independently
generated synthetic fixtures while preserving that omission.
\citep{exec24_unlearning_original,exec24_unlearning_software}

The first experiment should compare answer filtering, store deletion, cache
invalidation and model-side interventions under the same synthetic entity and
access policy. It should include a corpus-wide observer, split-value
reconstruction, retention of authorized tasks, correct refusal, raw-trace health,
fixed-scorer replay and a deliberately changed checkpoint. Public fixtures make
these contracts inspectable before costly bank trials. Success would justify
a bounded claim about declared surfaces and populations; universal forgetting,
legal compliance, causal validity and bank adoption remain separate questions.



% SOURCE docs/release/2026-10-29/source/docs/writing-repair/2026-10-29/src/bc9add2c0e3aeb85.tex CHAR 0
\section{ChronoLLM: domain-specific code assistance and the limits of simulation evidence}
\label{sec:exec24-chronollm}
A bank turning a collections question into executable analysis faces two
different tasks. It must express the intended decision, and it must implement
that decision in a changing software environment. Domain-specific code
assistance can reduce the second burden without solving the first. Wang and
coauthors study this distinction in an adjacent engineering setting:
ChronoLLM customizes language models to generate PyChrono scripts for
mechanical simulation. The inspected original is the 19 August 2025 arXiv
preprint, not the later journal version identified by the discovery index.
It supplies a useful design analogue for causal tools, rather than evidence
that its models identify bank interventions.~\citep{exec24_chronollm_original}

\subsection{Version-bound knowledge and expert-curated training}
The source describes stale and mixed API knowledge, with examples of renamed
vector, gravity, collision and visualization interfaces. Its proposed training
pipeline combines continued pretraining on domain material with instruction
fine-tuning. The instruction collections include simulation generation,
natural-language-to-API mappings, API explanations and reasoning examples.
Expert contributors clean scripts and comments; model-generated descriptions
and transformations supplement that work. The paper also describes filtering,
privacy checks, deduplication and final quality review. These are reported
curation procedures, not evidence that we acquired or independently audited
the full training corpus.~\citep{exec24_chronollm_original}

For the proposed causal library, this suggests binding every example to an
exact package version, method identifier and input/output contract. A
natural-language request for a default-risk comparison should not silently
select an obsolete method, change time zero, collapse competing events, or
replace a requested effect with a prediction. Retrieval and small curated
examples are plausible first prototypes; fine-tuning is a conditional
alternative when repeated version-specific failures justify its preparation
and maintenance cost. Porting the whole simulator or treating its generated
code as a causal engine would answer a different question.

The paper distinguishes in-context learning, which supplies examples without
changing model weights, from supervised fine-tuning and parameter-efficient
adaptation. Its broad architecture and model-selection discussion is not
adopted as a universal ranking. Our comparison would hold the selected
backend revision, prompt, package/API version, task family and allowed repair
budget fixed before estimating an improvement. An API spelling correction
and a correction to the scientific target must be separately observable.
\citep{exec24_chronollm_original}

\subsection{A score improvement depends on what the score measures}
The original distinguishes code similarity, execution tests and a separate
LLM judge. Its Table~7 reports similarity and judge scores under reference
code, documentation, or their combination. These are different measurements.
For the reported GPT-4o-mini profile, the combined reference/documentation
judge score rises from 41.80 to 68.03 after supervised fine-tuning, a difference
of 26.23 score units. The documentation-only score falls from 43.22 to 38.03,
a difference of $-5.19$. For the reported 8B profile, the combined judge score
is 39.06 with in-context learning and 38.78 with fine-tuning, a difference
of $-0.28$. Table~8 also contains lower judge scores for the 70B LoRA profile
than for its pretrained comparator. The exact subtractions reproduce the
displayed table entries; they are neither new experiments nor significance
tests.~\citep{exec24_chronollm_original,exec24_chronollm_reference}

Thus the paper's broad improvement language needs metric-specific reading.
Similarity to reference code does not by itself establish functional or
physical correctness. A judge score is not a probability that the output
passes a causal identification contract. The text discusses pass-at-$k$ and
attributes judge-correlation findings to a companion study, but the inspected
Tables~6--8 do not supply a numerical pass-at-$k$ series. The conclusion's
mention of a Hacked SimBench baseline is not accompanied there by an
identified table row or reproducible comparison definition. These gaps call
for the companion original, exact task/split manifests, scorer revision and
run artifacts before borrowing a stronger empirical claim. They do not
justify silently inventing missing results.~\citep{exec24_chronollm_original}

The two case studies show generated double-pendulum and vehicle scripts.
The authors acknowledge that their tasks may be in distribution relative to
training and describe the scripts as starting points for expert improvement.
The double-pendulum discussion distinguishes API errors from incorrect world
construction, axes and units. This is especially relevant to causal analysis:
a script can run while implementing the wrong population, assignment law or
intervention. The source reports example execution, but no acquired script,
simulator or model is executed in this proposal review.
\citep{exec24_chronollm_original}

\subsection{Executable worlds need not identify the intervention}
Consider an authored banking example with binary latent type $U$, equally
likely to be zero or one, and an observed action $A=U$. Let $Y$ denote
incremental account revenue in dollars. Two candidate structural worlds are
\begin{equation}
  \mathcal M_1:\ Y=10A+20U,
  \qquad
  \mathcal M_2:\ Y=30U.
  \label{eq:exec24-chronollm-worlds}
\end{equation}
Both generate the same observed pairs: $(A,Y)=(0,0)$ or $(1,30)$, each with
probability one half. The observed contrast is therefore USD~30 in either
world. Under an intervention that sets $A$ while preserving the distribution
of $U$, however, the average effect is USD~10 in $\mathcal M_1$ and zero in
$\mathcal M_2$. With a stipulated intervention cost of USD~5, the incremental
net value is USD~5 or USD~$-5$. The same observed data support opposite
decisions under these two structural explanations.
\citep{exec24_chronollm_reference}

Each world can be implemented with perfectly valid code and pass tests of its
own algebra. Execution cannot select which world describes the bank.
Conditional treatment overlap also fails here because action is fixed by
type. A useful agent should expose the unresolved assumption and request
appropriate additional information or a justified design. Public or generated
data remain valuable for testing whether the implementation preserves a
declared world and returns the required refusal; they do not turn that world
into an empirically verified bank mechanism. This provides a small exact
oracle for the proposed business-question-to-code workflow without paying
for an authentic bank experiment first.

\subsection{Full costs, repair accounting and conditional reuse}
The paper describes a small fine-tuning example taking roughly twenty minutes
and costing roughly USD~3 after the time-consuming data preparation was
complete. This is a reported historical example, not a current API quote,
full preparation cost, or measured bank productivity gain. Its model and task
profile must remain attached to the claim.~\citep{exec24_chronollm_original}

For an independent stipulated accounting witness, combine an assumed USD~3
API charge with thirty minutes of expert work at USD~60 per hour. Total cost
is USD~33. Suppose eight requested tasks yield six compilable scripts, four
semantically valid simulations and two causally acceptable outputs. Cost is
USD~4.125 per request, USD~5.50 per compilable script and USD~16.50 per
accepted output. These invented inputs show why expert repair, refusal and
failure denominators must be retained; they are not a reconstruction of the
paper's experiment.~\citep{exec24_chronollm_reference}

The proposed route is therefore reference use of its version-aware curation
and evaluation ideas, with an independently designed causal test suite.
The method owner defines intervention, units, clocks and admissible refusal;
the software owner pins package and interface behavior; the evaluation owner
holds out task families and versions, records all attempts and expert edits,
and compares retrieval, examples and fine-tuning under matched costs.
Unlearning stale APIs and richer tool or multi-agent interaction are future
directions in this preprint, not demonstrated universal remedies. Its
supplementary and declaration sections also use different availability
language: models/API access on request, a reported SimBench repository, and
a promised final code/data link. Exact artifacts, rights and reproducibility
therefore remain conditions for adoption.~\citep{exec24_chronollm_original}

This adjacent study supports a concrete public-prototype design while leaving
scientific identification, native adoption and bank value unestablished.


% SOURCE docs/release/2026-10-29/source/docs/writing-repair/2026-10-29/src/732e9f0b1fa8e4a7.tex CHAR 0
\section{Prompt optimization must preserve the question, evidence and authority}
\label{sec:exec24-prompt-engineering}

An analyst can improve a prompt while making a banking answer less defensible.
More detailed instructions may introduce an unstated population assumption;
a supplied worked solution may make a benchmark easier without improving
independent problem solving; a tool-using agent may execute a valid command for
an invalid causal question. The relevant product question is therefore which
prompting design improves an explicitly defined service under a fixed evidence
and authority contract. Chen and colleagues' \emph{Unleashing the potential of
prompt engineering for large language models} provides a broad taxonomy of
instructions, demonstrations, decomposition, repeated sampling, search,
optimization, multimodal prompting, evaluation and security concerns.
\citep{exec24_promptengineering_original}

The inspected source is the 58-page arXiv revision 6 dated 11 May 2025, the
latest revision listed before the client's 7 October 2026 cutoff. The initial
23 October 2023 version has 25 pages and a different title; it is not silently
substituted for revision 6. The archive also supplies a journal identity, but
the final journal bytes have not been equated with this preprint. We use the
survey as a primary record of its authors' organization, examples and caveats.
Its references do not themselves qualify the methods, packages or empirical
rankings it describes. No model is run in this review.

\subsection{Extra context can change the target instead of identifying it}

The source's introductory examples compare a short question with more elaborate
prompts and generated context. In its entrance-width example, later answers
consider height, terrain, weight and access restrictions that were not fixed in
the question. Those may be useful questions to ask, but neither physical
possibility nor an expanded list of possibilities defines a probability law over
them. A preferred answer based on presumed prevalence imports an assumption.
The source describes richer answers as more coherent; our narrower conclusion
is that coherence must be evaluated separately from support by the given
information. Figure text was inspected, without fresh visual review.
\citep{exec24_promptengineering_original}

For a banking analogue, let $d$ be the reduction in default probability from an
intervention at a specified horizon, $L=1000$ dollars the loss per default and
$c=20$ dollars the action cost. Under the stipulated linear loss calculation,
net value is $V(d)=Ld-c$. Suppose the available information only restricts
$d\in[-0.10,0.10]$. The compatible values are
\[
 V(d)\in[-120,80]\quad\text{dollars per eligible account},
 \qquad V(d)>0\ \Longleftrightarrow\ d>0.02.
\]
If an elaborate prompt inserts $d=0.05$ as plausible generated knowledge, it
returns USD 30. The arithmetic is correct conditional on that value; the
supplied information did not identify it. The appropriate output retains the
range, break-even condition and the need for additional evidence. A user can
explicitly request the five-percentage-point scenario, but a scenario must be
labeled and must not overwrite the identified-set result.
\citep{exec24_promptengineering_reference}

The translation interface should therefore return missing-premise questions
alongside its candidate target and graph. It should distinguish a quoted fact,
an analyst-supplied restriction, an assumed scenario and model-generated
context. Retrieval adds a document, not automatic truth; the document's
population, time and authority still require interpretation. The proposal
prefers a structured questionnaire when that is sufficient. Generated knowledge
is suitable for proposing hypotheses and sensitivity cases, subject to explicit
status, rather than filling causal evidence gaps invisibly.

\subsection{A task graph is not a causal graph}

The survey discusses chains, trees and graphs of generated intermediate
information, and a decomposed workflow with a controller and subtask handlers.
These are descriptions of computation. An edge from a retrieved statement to
a proposed answer records a dependency of the answer on text; it does not
assert that intervening on an economic variable changes another economic
variable. The graph used to schedule computation and the structural graph used
for identification consequently need different types, identifiers and
validation rules. Converting one into the other requires a scientific argument.
\citep{exec24_promptengineering_original}

This distinction guides integration. A deterministic handler can parse a
versioned variable dictionary; an approved estimator can consume a reviewed
query; a model can propose an explanatory draft. Each handler should declare its
inputs, outputs, preconditions, failure statuses and artifact versions. The
controller must preserve an unidentified or refused status instead of translating
it into a numerical default. Repeating a subtask must create a traceable attempt,
not erase the earlier refusal. These are proposed CLI/MCP contracts, not a claim
that a native agent framework already enforces them.

An authored reliability example shows why local component scores do not settle
the end-to-end claim. Let $E_i$ mean that stage $i$ is correct, for five stages,
and stipulate $P(E_i)=19/20$ for every stage. If their correctness events are
independent, complete correctness has probability
\[
 P\!\left(\bigcap_{i=1}^{5}E_i\right)
     =\left(\frac{19}{20}\right)^5
     =\frac{2476099}{3200000}\simeq0.773781.
\]
With only those marginal probabilities known, the union bound and event
containment instead give
\[
 0.75\leq P\!\left(\bigcap_{i=1}^{5}E_i\right)\leq0.95.
\]
Disjoint five-percent failure events attain the lower bound; identical failure
events attain the upper bound. Independence is additional information, not a
default consequence of modular code. These are unconditional event assumptions
under a declared task law, not measurements of a five-stage application.
\citep{exec24_promptengineering_reference}

For sequential evaluation, later stages may receive wrong earlier outputs or
may never run after a refusal. Report both attempted-stage denominators and
complete-task outcomes under the real controller policy. In particular,
testing each handler only on correct supplied inputs establishes a different
property from testing the composed workflow on manager requests.

\subsection{Prompt selection is learning and needs held-out evaluation}

The survey covers selecting uncertain questions for annotation and optimizing
prompts using feedback, search and learned rewriting. It also reports that
different tasks, models and metrics limit comparisons, and that automated
prompting can fail to improve on simple baselines. These are source-level
observations, not reproduced rankings. An optimized prompt is fitted to
information: developer trials and expert edits are part of training even when
model weights stay fixed.~\citep{exec24_promptengineering_original}

Consider a constructed null example with $K=20$ candidate prompts and $n=4$
development tasks. Every prompt has true success probability $1/2$. Assume
success indicators are independent across prompts and tasks. For prompt $j$,
let $S_j$ be its development successes. Then $S_j\sim\operatorname{Binomial}(4,1/2)$.
Selecting the largest score gives
\[
 P\{\max_j S_j=4\}=1-(15/16)^{20}\simeq0.724941.
\]
Writing $F(m)=P(S_j\leq m)$, the expected best development accuracy is
\[
 E\!\left[\frac{\max_j S_j}{4}\right]
 =\frac14\sum_{m=1}^{4}\{1-F(m-1)^{20}\}
 \simeq0.931096.
\]
Nevertheless, a fresh test independent of selection still has expected
accuracy $1/2$ for the selected prompt. A development winner above 93 percent
can arise in this deliberately small, heavily searched example without any
improvement in true performance. The numerical probabilities depend on the
declared independence law; they are not universal estimates of prompt overfitting.
\citep{exec24_promptengineering_reference}

The banking protocol should therefore reserve unseen tasks or task families,
record the complete search and repair budget, and separate development examples
from scoring evidence. Frozen model and retrieval versions make a comparison
interpretable but do not themselves guarantee generalization. Where tasks share
a template, customer history or policy source, that clustering affects the unit
of evaluation. Repeated paraphrases of one causal problem must not be counted
as independent scientific problems without justification.

The survey's discussion of supplied solution chains is particularly relevant:
it acknowledges that supplying ground-truth intermediate solutions limits what
a high solve rate demonstrates. Our evaluation should label an oracle-assisted
condition separately from unaided translation or identification. Likewise,
zero-shot describes the declared examples in a prompt, not verified absence of
the problem from a model's pretraining data. Neither a polished explanation nor
agreement among human judges independently establishes causal correctness.
\citep{exec24_promptengineering_original}

\subsection{Prompt text cannot grant tool authority}

The survey includes adversarial inputs, poisoned data, backdoors, prompt injection
and leakage. Some of its language is stronger than a bounded experiment can
warrant. This proposal retains the threat categories without adopting a claim
that prompt design ensures secure deployment. Full source inspection of an
attack or defense, its version and threat model would precede implementation.
\citep{exec24_promptengineering_original}

A proposed tool gate separates three logical conditions: $T$, a trusted request
authorizes the action; $A$, the action is allowed by the relevant policy; and
$H$, the request's artifact identities match those reviewed. Define admission
as $T\land A\land H$. Enumerating the eight Boolean configurations admits one
and rejects seven. This exact truth table demonstrates only the declared
predicate. It does not prove authentication, policy completeness, correct hash
binding, parser safety or atomic execution in a real system.
\citep{exec24_promptengineering_reference}

For example, retrieved text may recommend replacing a reviewed causal graph.
That text is evidence to inspect, not an authorized request to overwrite it.
A proposed MCP tool should return a candidate revision with source provenance;
an effect-producing action needs the separate trusted request and relevant
policy checks. The same distinction applies to CLI flags and shell commands.
Read-only retrieval, query construction and execution should have explicit
operation types, with the relevant authority verified at each external action.
Textual role instructions are useful guidance but are not the authority record.

\subsection{Conditional adoption and complete economic comparisons}

The choices are not simply an advanced prompt or no assistance. The proposal
compares a deterministic form, an expert-written template, retrieved examples,
bounded decomposition and adaptive optimization on the same declared tasks.
Use a simple template when it meets the target; add retrieval when source
coverage and provenance justify it; use original, qualified framework interfaces
when reusable orchestration earns its cost. A whole-framework port or compulsory
agent layer needs a demonstrated requirement. The survey supplies discovery
leads for these alternatives, not software rights or functioning integrations.

In one authored cost allocation, ten attempted tasks cost USD 4 in model calls,
USD 2 in retrieval and USD 24 in review. If three produce accepted outputs,
cost is USD 3 per attempt but USD 10 per accepted output. Actual pricing,
latency, maintenance, failures and correct-refusal value remain to be measured.
A cheap request is not automatically a cheap decision. Public documents and
constructed causal worlds permit this prototype comparison before authentic
bank experiments, while source rights, representation and transfer limits remain
explicit.~\citep{exec24_promptengineering_reference}

The resulting adoption rule is conditional: keep a prompt or agent design only
when its declared target, evidence handling, complete-task behavior and cost
justify it against the simpler alternative. This source broadens the taxonomy
and sharpens test requirements; it leaves native integration, complete frontier
coverage and bank acceptance unestablished.


% SOURCE docs/release/2026-10-29/source/docs/writing-repair/2026-10-29/src/f5ead26280448b86.tex CHAR 0
\section{From scientific assistance to a testable banking research loop}
\label{sec:exec24-scientific-method}

An automated research loop needs a scientific question that survives every transformation.
A servicing analyst may ask whether an earlier contact changes ten-day cash receipts,
while a text assistant extracts an eventual repayment flag and an optimizer maximizes
that flag. Each component can operate correctly while the composition answers a
different question. This is the practical value of examining the whole scientific
cycle, rather than treating a better hypothesis generator as a complete causal system.
Zhang and colleagues organize their perspective around observation, hypothesis,
experimentation and automation, with human scientific goals and evaluation as essential
constraints. We use their May 2025 preprint as a primary source for that argument
and its caveats. Its descriptions of other systems are leads to their originals,
not independent qualification of those systems.~\citep{exec24_scimethod_original}

The source is the 45-page version of \emph{Advancing the Scientific Method with Large
Language Models: From Hypothesis to Discovery}, arXiv 2505.16477v1, dated 22 May 2025.
A later journal record has a different title and the same thirteen authors. The
journal full text was not available in this inspection; neither exact equivalence
nor journal-specific claims are assumed. This distinction matters because a review
may change its qualifications or examples between versions. The preprint's account
supports a design vocabulary and criticism, not an empirical estimate of bank
productivity or a theorem that autonomous discovery will succeed.

\subsection{Observation is part of the estimand}

The perspective treats observation as a reduction of complex information, including
annotation, classification and extraction. It also acknowledges information loss when
domain data are converted into a representation accepted by a model. In banking the
representation is consequential: a loss of timing, cohort membership or assignment
information can make an otherwise answerable intervention question unanswerable.
Let $H$ denote the original account history and let $Z=g(H)$ be the assistant's
extracted representation. For a target $q(H)$, a deterministic reconstruction from
$Z$ exists on the declared history class only if $q$ is constant on every fibre
of $g$. Indeed, if $g(h_1)=g(h_2)$ but $q(h_1)\ne q(h_2)$, any proposed function
$f$ would have to give two values to $f(g(h_1))$. Conversely, constancy allows
$f(z)$ to be defined by the common target value on its fibre. This elementary
statement concerns information preservation, not stochastic causal identification.

Consider two stipulated histories of a USD 10,000 obligation: full repayment on
day 5 and full repayment on day 25. The extraction of eventual repayment is one
in both histories, while the ten-day cash target is USD 10,000 in the first and
zero in the second. An accurate eventual-repayment classifier therefore does not
resolve the liquidity target. Keeping raw histories, event times, availability
times and extraction provenance is a substantive alternative to committing early
to a convenient label. It costs more storage and review effort, but allows later
questions to be assessed for information sufficiency. Neither retained detail nor
a high extraction score alone establishes truthful assignment or absence of
unmeasured confounding.~\citep{exec24_scimethod_reference}

A business-to-causal translation should consequently produce a structured question:
population, intervention and comparator, outcome, horizon, unit, information available
at decision time, and loss function. A generated diagram then proposes mechanisms
consistent with those definitions. Its edges need separate provenance: a substantive
restriction, a temporal impossibility, an inspected empirical result, or an
unresolved hypothesis. Textual plausibility cannot decide between observationally
equivalent mechanisms. Where the question changes after seeing results, that change
belongs in a new target and evaluation record, rather than silently changing the
meaning of the original request.

\subsection{Hypothesis generation and verification have different jobs}

The source explicitly discusses a possible creative role for speculative generation,
while requiring reliability in experimental work. This is useful only when an
unsupported suggestion remains visibly hypothetical. A proposed edge from contact
to repayment delay is a candidate for comparison with alternatives,
not a fact to insert into every downstream adjustment set. Competing stories may
include selection by latent repayment capacity, contact triggered by early warning
signals, and a genuine contact effect. Their distinguishing implications must be
derived before choosing an observation or intervention that could separate them.
Retrieval can improve access to documented premises; it cannot turn a similar
population, treatment or observation clock into the same one.

The perspective's discussion of formal verification contains an especially important
qualification: translating a generated hypothesis into a formal language can itself
introduce error. For example, the goal of cost strictly below USD 100 is
$c<100$. An assistant translating it as $c\leq100$ may produce a perfectly verified
solution at $c=100$ that violates the original goal. A proof checker is entitled to
accept the weaker formal specification. The missing obligation is correspondence
between business intent and formal premises. Our exact boundary witness does not
test a proof assistant; it explains why a verified artifact must bind the agreed
question, units, inequalities, domain and assumptions as well as the solver output.
\citep{exec24_scimethod_original,exec24_scimethod_reference}

There are therefore several distinct checks. A parser can establish syntactic
validity; a type checker can reject a unit or shape mismatch; an exact small oracle
can test a proposed algorithm on a declared family; a proof can establish a theorem
under explicit premises; and an experiment can challenge those premises in a stated
population. Failure at any layer is informative, and passing one layer does not
eliminate the others. Generated executable code can falsify a claim on a witness
without proving it for all permitted inputs. A successful formal derivation also
does not make a bank's observational assignment mechanism exogenous.

\subsection{Repeated agreement is not calibrated scientific confidence}

The preprint surveys repetition, self-correction, debate and external feedback, but
also reports limitations of reasoning, planning and self-explanation. It calls for
system-level trust assessment that includes tasks, tools, workflows and feedback.
Its proposed analogy with statistical confidence is a research direction; it does
not supply a generally valid confidence interval for a causal answer produced by
an arbitrary agent. We should retain that question rather than converting it into
a claimed product feature.~\citep{exec24_scimethod_original}

An exact calculation illustrates the issue. Suppose five judges each have error
probability $p=0.1$ on a stipulated request law. Under independent errors the
majority-error probability is
\[
 \sum_{k=3}^{5}\binom{5}{k}p^k(1-p)^{5-k}
 =\frac{107}{12500}=0.00856.
\]
If every judge instead shares the same Bernoulli error, each still has marginal
error $0.1$, but majority error remains $0.1$. The shared failure could arise from
the same incorrect source, an identical omitted confounder or a common mistranslation.
The example does not estimate any real model's accuracy, nor does it establish a
universal worst case over all dependence structures. It shows that marginal
performance and the number of agreeing judges do not determine the error of their
composition. A new model name or a different prompt is insufficient evidence of
independence.~\citep{exec24_scimethod_reference}

The bank-facing output should distinguish process completion, answer acceptance,
correct refusal, and unresolved cases. A confidence statement additionally needs
a defined event, population and task distribution, a sampling design, and a
justified construction evaluated on tasks not used for search or repair. An
independent causal oracle for generated models is one useful early tool. It can
check whether an agent preserves ambiguity or refuses an unidentified request,
instead of rewarding a fluent numerical answer. Evaluation on public observations
without intervention truth supports more limited checks of provenance,
reproducibility and assumptions. A model's own explanation can guide debugging,
but the perspective's concerns about self-explanation prevent treating it as a
faithful account of the model's internal mechanism.

\subsection{What public evidence can establish before bank experiments}

The product proposal should earn the expense of bank experiments through a staged
public and generated evidence portfolio. Generated structural models can supply
known intervention answers, nonidentification pairs, deliberate clock errors and
controlled measurement loss. Public empirical data can exercise extraction,
missingness, computation and reproducibility on realistic structures, with the
available causal truth stated separately. Open scientific tasks can test planning
and tool use, but their endpoints and safety constraints need not match banking.
This division lets inexpensive evidence test a concrete architectural proposition
without claiming that public data contain unavailable bank truth.

Transport remains a separate proposition. In a stipulated public environment,
contact and waiting yield net values USD 8 and USD 3 per account. A system can
select contact perfectly. In a second, explicitly different bank environment,
their values are USD $-4$ and USD 3. Applying the public choice incurs USD 7
regret relative to waiting. These are constructed utilities, not estimates. The
example isolates the need for a transport argument: good evaluation inside one
world does not establish which world the bank inhabits. Shared mechanisms,
covariates and measurement definitions may support a restricted bridge, but
those restrictions require evidence.~\citep{exec24_scimethod_reference}

For a first prototype, compare a fixed human-authored workflow, a retrieval-assisted
assistant, and a bounded planning assistant on the same frozen tasks and evidence.
The comparison should fix the allowed tools, retry budget, development/evaluation
split, answer and refusal rules, and costs before the final run. Include failures
to recover a paper, incorrectly interpreted clocks, ambiguous diagrams, unsupported
identification, malformed tool responses and stale results. A workflow that
correctly refuses an unanswerable request has a different achievement from one
that returns an accurate identified effect; both should be reported explicitly.
This is a proposed protocol, not an executed experiment or a claim of independent
acceptance. The number and mix of trials should follow the declared precision
and decision requirements, rather than an arbitrary success threshold.

\subsection{An auditable loop through the CLI and MCP}

The scientific loop motivates a concrete proposed interface. A question artifact
should bind the target, source identities and information clock. A diagram artifact
should carry alternative graphs and the provenance of constraints. An identification
request should return an estimand, an identified set, or a reasoned refusal. An
experiment proposal should bind the competing hypotheses and distinguishing
observations; an evaluation result should bind the exact task, methods and outputs.
CLI commands and MCP tools can expose these same typed artifacts, so that a
planner's change of wording cannot silently change the scientific target.

Acquisition, extraction, proof or numerical checking, and any execution have
different permissions and evidence duties. The proposal should keep hypothesis
content separate from permission to act. Before a tool invocation the host needs
to validate the requested operation, artifact identity and allowed context;
afterwards it needs to preserve the actual result and failure state. Retries
must retain earlier errors and cost, and resumed sessions must bind the same
frozen request. A local exact reference, an external package worker, and a
human-reviewed result are different execution modes. These contracts describe
required behavior; this source review has not implemented or qualified them.

Directly adopting a perspective as executable infrastructure is inappropriate:
it supplies an organizing argument, not a tested interface or algorithm. A useful
adoption is its separation of observation, hypothesis and experiment, accompanied
by the stricter causal and evidential contracts developed here. Porting an agent
framework is not justified by its appearance in the survey. The relevant package,
version, licensing, dependencies, termination and data boundaries need their own
primary review. A small explicit state machine remains a substantive alternative
to a general planner when the workflow is known; a planner may offer flexibility
when the next useful observation is genuinely unresolved. That flexibility must
be evaluated against its additional error, cost and review burden.

Finally, assume a fully counted exploratory workflow costs USD 240 for 40 attempts,
of which 30 complete their tools and eight deliver accepted scientific answers.
Its costs are USD 6 per attempt, USD 8 per tool completion and USD 30 per accepted
answer. Six correct refusals are reported separately; they are not silently added
to the answer denominator. All amounts and counts are stipulated. Recording
failed hypotheses and unsuccessful tool calls, consistent with the perspective's
case for learning from negative results, is necessary to estimate eventual
economics honestly. The contribution here is a testable proposal with explicit
information, translation, dependence and transport limits; native performance,
general algorithmic confidence and bank value remain unestablished.
\citep{exec24_scimethod_original,exec24_scimethod_reference}


% SOURCE docs/release/2026-10-29/source/docs/writing-repair/2026-10-29/src/60fc8d0b95950571.tex CHAR 0
\section{Causal affordances and executable tests for an agent-assisted library}
\label{sec:exec24-affordance-fuzzing}
A proposed causal assistant must select useful analytical actions and expose mistakes in the software that carries them out. These are related but distinct responsibilities. Reasoning about which familiar object can serve an unfamiliar purpose offers a controlled test of functional selection; generating programs that expose computational inconsistencies offers a test of software behavior. Neither alone establishes a banking intervention effect. The affordance study by Gjerde, Cheung and Lagnado and the TitanFuzz study by Deng and colleagues motivate a joint design in which agents propose actions and test inputs while explicit scientific and computational contracts determine what counts as success.~\citep{exec24_affordance_original,exec24_fuzzing_original,exec24_affordance_fuzzing_reference}

\subsection{What an affordance benchmark measures}
The inspected affordance text is the 21-page February 2025 author preprint. It studies 20 tasks in which an ordinary tool is unavailable and the respondent selects a substitute. The four-choice design includes objects associated with the ordinary tool but unsuitable for the task, a less associated suitable object and a distractor. Experiment one recruits 100 UK participants and compares their responses with historical model configurations under different prompts and temperatures. The paper explicitly acknowledges that repeated stochastic model runs are not independent human participants, although run-level averages are used as comparable data points.~\citep{exec24_affordance_original}

This is a useful challenge to association-based selection: the familiar object can be the wrong answer. It does not identify an external intervention effect from observational data. Task authors have specified the intended function and scored the offered alternatives; a bank still needs evidence about legal eligibility, timing, collateral, assignment and outcome mechanisms before translating a plausible suggestion into an action. The study's qualitative examples illuminate response patterns, but they were selected for notable differences and are not a separate representative sample of reasoning failures.~\citep{exec24_affordance_original,exec24_affordance_fuzzing_reference}

Experiment two retains 296 of 300 recruited UK participants and introduces nine-option distractor and image conditions. Its ordinary prompt permits longer answers than the first experiment's restricted prompt. The authors explicitly discuss this change as a possible explanation of part of the earlier prompting effect. The separate-option reasoning prompt is not evaluated in the nine-option condition because the resulting answers would be long. This design does not support a fully crossed estimate of every prompt-by-difficulty interaction.~\citep{exec24_affordance_original}

The reported results include configuration-specific similarities to human accuracy, performance declines with additional distractors, and different responses to images across models. A nonsignificant difference is not an equivalence result: equivalence requires a justified margin and an appropriate test. The paper also distinguishes newly written tasks from structurally new problems and acknowledges incomplete control over training-data overlap. Its behavioral observations do not determine the internal mechanism or demonstrate real-world tool use. Historical model names describe the studied conditions, not current product specifications or a claim of superiority today.~\citep{exec24_affordance_original,exec24_affordance_fuzzing_reference}

\subsection{Translate functional selection into a banking contract}
Consider a stipulated liquidity problem: obtain eligible funds within one business day. An action's required properties are a settlement delay no greater than one day, eligible collateral and an existing authorization. Four authored choices are a familiar facility settling in two days, an unfamiliar facility satisfying all three conditions, an immediate facility without eligible collateral, and an immediate facility lacking authorization. The unfamiliar facility is the unique feasible choice. Adding five infeasible distractors leaves that answer unchanged. These are declared fixture properties, not claims about actual facilities or law.~\citep{exec24_affordance_fuzzing_reference}

The assistant can decompose the request into constraints and attach source passages, dates and unresolved interpretations to each property. It should then distinguish feasibility from preference. If two choices satisfy the constraints, ranking them requires costs, risk and the bank's objective; if collateral eligibility is unknown, the answer must retain that uncertainty. A familiar label or an explanation that sounds financially plausible cannot supply the missing fact. In graph construction, the analogous step is to distinguish candidate variables and functional mechanisms from assumptions needed for causal identification.~\citep{exec24_affordance_fuzzing_reference}

Test representation and substantive meaning separately. Renaming facilities or permuting choices should preserve this fixture's answer; changing settlement from one to two days should not. A unit conversion must preserve the same economic threshold. Changing an observation timestamp can invalidate a proposed covariate. These transformations create an auditable test family whose required behavior is defined before a model generates an answer. They do not establish that the chosen public task distribution represents a bank's future decisions.~\citep{exec24_affordance_fuzzing_reference}

Choice count also changes a naive accuracy baseline. In an authored comparison, accuracy 0.85 among four choices has chance-normalized score
\[
 (0.85-1/4)/(1-1/4)=0.80,
\]
whereas accuracy 0.75 among nine choices gives
\[
 (0.75-1/9)/(1-1/9)=23/32=0.71875.
\]
Both raw and normalized scores are informative; neither removes changes in task difficulty or establishes causal competence. The numbers are illustrative and are not the paper's reported contrast.~\citep{exec24_affordance_fuzzing_reference}

Repeated outputs do not automatically increase the number of independent problems. Suppose there are 20 independent task clusters, each with 25 unit-variance errors, equicorrelated at a stipulated 0.60 within task. The mean over 500 outputs has variance
\[
 \frac{1+(25-1)0.60}{20\cdot25}=0.0308,
\]
compared with 0.002 under independence. The equivalent independent count for this variance calculation is $2500/77$, about 32.47. A different dependence law gives a different answer. This witness motivates retaining task identity, shared context, prompt selection and model revision in the evaluation design. It is not an estimate of the study's correlation or a reanalysis of its significance tests.~\citep{exec24_affordance_fuzzing_reference}

\subsection{Generate difficult programs while preserving the test oracle}
TitanFuzz supplies a complementary engineering idea. Its inspected March 2023 author version combines a generative model for seed programs with an infilling model for mutations. Mutation operators alter arguments, surrounding code and methods; adaptive operator selection rewards valid generations, while a structural fitness score favors dataflow depth and unique API calls and penalizes repetition. The paper uses successful execution with target invocation as a program-validity criterion. Validity, diverse execution and discovery of a real defect are separate outcomes.~\citep{exec24_fuzzing_original}

The paper's principal oracles compare CPU and GPU computations and identify crashes. Comparison tolerances address ordinary numerical variation. The experiments cover historical PyTorch and TensorFlow configurations, with a one-minute per-API budget for the main comparison and a four-minute budget for bug finding. Reported line coverage concerns the measured Python code; it does not establish that every lower-level path, relevant banking contract or possible defect was exercised. The paper reports 65 candidate findings, 53 developer confirmations and 41 previously unknown bugs at its reporting date, with eight rejected findings. One rejected example has intentionally implementation-dependent behavior. This distinction matters: a discrepancy is evidence to investigate, not automatically a valid defect.~\citep{exec24_fuzzing_original}

The inspected repository revision is dated September 2023 and reports later totals of 55 confirmations and 44 previously unknown bugs. These dated source statements remain separate; we have not independently adjudicated the individual reports. The revision's complete tree has 88 entries. Its README identifies the research paper, describes supplied historical seed artifacts and points to a seed-generation entrypoint absent from that tree. The external seed archive has not been acquired in this review. Thus the repository does not by itself establish complete reproduction of the paper's generation process.~\citep{exec24_titanfuzz_source}

A version contract is also necessary. The inspected requirements leave several major dependencies unpinned, while the Docker recipe names TensorFlow 2.11.0rc1, differing from the paper's TensorFlow 2.10 experiment. The source includes process recovery, timeout and skip statuses, but these mechanisms alone do not establish containment, reproducible numerical behavior or target-platform acceptance. No native model, downloaded program, container or package is executed in this review.~\citep{exec24_titanfuzz_source}

\subsection{A numerical agreement check is not an economic acceptance rule}
The inspected comparator declares relative tolerance 0.01 and absolute tolerance 0.001 and calls library closeness functions with equal-NaN handling. Some unhandled value types fall through to a true return. These are static source observations whose consequences need carefully specified native tests. They should not be converted into a claim that a particular version has been observed to misclassify a banking result.~\citep{exec24_titanfuzz_source}

An authored scalar example illustrates the required distinction. Define the comparison rule
\[
 |x-y|\leq a+r|y|,
\]
with $a=0.001$ and $r=0.01$. If two backends return 1,000,000 USD and 1,005,000 USD, their difference is 5,000 USD and the tolerance is 10,050.001 USD. This rule accepts agreement even though the difference exceeds a separately stipulated 1,000 USD business tolerance. The scalar rule is our explicit witness using the inspected constants; it is not a native execution of either numerical library.~\citep{exec24_affordance_fuzzing_reference}

Absolute tolerances carry units. Comparing zero with 0.0005 dollars passes an absolute threshold of 0.001 dollars. Expressing the same values in cents while retaining the numerical threshold 0.001 fails; the threshold must also be converted. A causal API therefore needs units, estimand and error budgets in its artifact contract, including any conversion performed by a CLI or MCP consumer.~\citep{exec24_affordance_fuzzing_reference}

Backend agreement can also preserve a shared error. If the exact truth is 100, both backends return 110 and action occurs above 105, perfect agreement produces the wrong action. In inference, a shared identification mistake is even less accessible to numerical comparison. Use independent mathematical references, exact finite models, invariance tests and deliberately invalid requests alongside differential tests. Refusal on an unidentified effect is part of correct behavior, not a program-generation failure to discard.~\citep{exec24_affordance_fuzzing_reference}

\subsection{Adoption, interfaces and the public prototype}
The defensible proposal is to adopt the testing pattern conditionally while treating the available code as a separately qualified dependency. No license, copying or notice file appears in the inspected tree, and the conventional license request returns 404. This bounded observation leaves reuse rights unresolved; it does not establish that no permission exists elsewhere. Direct code adoption or porting should wait for an applicable grant and review of dependencies, model artifacts and supplied seeds. A project-owned harness can be designed from the inspected scientific ideas without declaring the downloaded implementation accepted.~\citep{exec24_titanfuzz_source,exec24_affordance_fuzzing_reference}

A public prototype should freeze a versioned task schema before testing: business objective, units, clocks, graph and identification contract, permitted operations, expected success or refusal, and the numerical tolerance. The agent proposes candidate programs or structured requests. A controlled runner records their source identity, environment, resource limits and statuses; an independently specified oracle evaluates outputs. The scientific checker must be allowed to reject a numerically successful calculation when the causal target is unsupported. Corpus revisions and model changes require a new held-out comparison.~\citep{exec24_affordance_fuzzing_reference}

For CLI and MCP consumers, generation, execution, adjudication and reporting should exchange separate immutable artifacts. A request identifies its oracle and authority scope; a response records all attempts, timeouts, duplicates, failures, unresolved discrepancies and accepted outputs. Sibling repositories are prospective consumers of this protocol, not demonstrated integrations. Before copying a discovered reproducer into another workspace, review its permissions, inputs, resource behavior and dependence on the environment. These are proposed product obligations, with no cross-repository deployment asserted.~\citep{exec24_affordance_fuzzing_reference}

An exact cost fixture starts with 100 attempts, 60 runnable programs and 20 duplicates among the runnable outputs, leaving 40 unique runnable programs. Ten candidate discrepancies comprise six confirmed defects, two intended differences and two unresolved cases. At a stipulated total cost of 600 USD, costs are 6 USD per attempt, 10 USD per runnable program, 15 USD per unique runnable program and 100 USD per confirmed defect. These denominators include different services. They are not market prices or a measured return on investment. A future comparison must include generation, hardware, containment, failed runs and human adjudication, then connect accepted outputs to decision value.~\citep{exec24_affordance_fuzzing_reference}

The combined lesson is practical. Functional decomposition can help formulate an analysis and adversarial program generation can expose implementation errors. The merit of a public experiment comes from its explicit semantics, independent oracle and complete accounting. Neither task accuracy nor software coverage supplies missing assignment information, resolves observationally equivalent causal worlds or substitutes for later institutional acceptance.~\citep{exec24_affordance_fuzzing_reference}


% SOURCE docs/release/2026-10-29/source/docs/writing-repair/2026-10-29/src/8e264b8f1a8ff636.tex CHAR 0
\section{Choosing an agent framework without changing the causal question}
\label{sec:exec24-framework-census}

An agent framework can help a bank turn a loosely specified business request
into a sequence of inspectable analytical tasks. It can retain a conversation,
route work, call a causal engine and recover an interrupted calculation. The
scientific question remains separate: what intervention is proposed, for whom,
at what information time, under which assumptions, and with what uncertainty?
The useful comparison is therefore between concrete execution contracts and
their ability to preserve that question. A returned sentence, a completed
workflow and a qualified causal recommendation are distinct outcomes
\cite{exec24_framework_contracts}.

Consider a request to increase retention among accounts at risk of leaving.
An assistant can elicit whether the action is a fee waiver, a contact policy
or a credit-limit change; identify the decision date; retrieve a data
dictionary; propose competing graphs; and submit an explicit intervention
query. Every step can be useful before a production product exists. Public
fixtures can test whether the assistant preserves treatment versions,
account histories, unresolved confounding and refusal messages. A later
bank experiment is justified only after these cheaper checks demonstrate
that the proposed chain addresses the relevant scientific problem
\cite{exec24_framework_contracts}.

\subsection{Source profiles and dependencies are part of the comparison}
The inspected LangChain repository contains distinct declarations for
modern LangChain 1.4.3, core 1.6.7 and classic 1.0.8. The modern package
declares a LangGraph range, rather than implementing every execution layer
itself. The separately retained LangGraph source declares 1.2.14.
Comparing the packages as independent causal engines would double-count
the shared orchestration dependency and misstate what either provides
\cite{exec24_framework_langchain,v12_la04_langgraph}.

The other inspected profiles are Semantic Kernel's Python 1.45.0,
DSPy 3.4.0, LlamaIndex core 0.14.25 and a Haystack source declaration of
3.4.0-rc0. A release-candidate label must survive the comparison; it is
not silently promoted to a stable distribution. LlamaIndex's workflow
modules re-export another package. The traced repository resolves to
\texttt{run-llama/llama-agents}, where the selected package declares
\texttt{llama-index-workflows} 2.25.0. This declaration fits core's
specified range, but no solved environment or native compatibility test
has been performed
\cite{exec24_framework_semantic_kernel,exec24_framework_dspy,
exec24_framework_llama_index,exec24_framework_haystack,
exec24_framework_workflows}.

These are authenticated source revisions with commit dates no later than
7 October 2026, inspected from retained original bytes. They are not
claims about the latest release, a bank's installed version, or all
permitted dependency combinations. Root notices for the selected
LangChain, Semantic Kernel, DSPy, LlamaIndex and workflow sources are MIT;
Haystack's selected root notice is Apache License 2.0. AutoGen's earlier
review keeps its AgentChat code notice separate from its documentation
notice. Complete dependency, model-weight, connector and hosted-service
rights remain component-specific inputs
\cite{exec24_framework_langchain,exec24_framework_semantic_kernel,
exec24_framework_dspy,exec24_framework_llama_index,
exec24_framework_haystack,exec24_framework_workflows,v12_la04_autogen}.

\subsection{LangChain and LangGraph: graph recovery is not side-effect recovery}
The modern LangChain factory constructs a state graph and passes a
checkpointer, store and interrupt boundaries into compilation. This is
a concrete route for retaining the business brief, candidate graphs,
assumption decisions and tool results. The factory's configured recursion
limit is a graph-execution control; it is not a token budget, monetary
budget or end-to-end service deadline
\cite{exec24_framework_langchain}.

The retained LangGraph interrupt documentation states that resumption
restarts the node and re-executes its logic. A checkpointer is required
for the documented interrupt route. These properties support recoverable
analysis, but preceding external actions need their own replay contract.
If a node creates a paid external job before the interrupt, restored graph
state does not by itself prove that the job can safely be created again
\cite{v12_la04_langgraph,exec24_framework_contracts}.

The inspected modern retry middleware makes a useful further distinction.
It counts the initial attempt plus the configured retries, propagates
graph control-flow exceptions, and treats nonmatching exceptions separately.
After exhaustion, the tool middleware can return a tool message explicitly
marked as an error; the model middleware can return an assistant message
containing failure information. Both also permit an error-raising route.
A wrapper that tests only whether a message exists would erase the
difference between recovered output and reported failure
\cite{exec24_framework_langchain}.

An authored banking example specifies the resulting public test. Suppose
an external service accepts a 500 USD action, but its acknowledgement is
lost. Three submissions without a shared deduplication key can represent
1,500 USD of actions. Under an explicitly stipulated service contract
that deduplicates the same key, they represent 500 USD. The 1,000 USD
difference is not an observed framework defect; it is the consequence
of different external contracts. The proposed adapter must preserve the
logical action key, attempt identity, acknowledgement state and recovery
query. Retrying an unidentified causal request also does not repair its
scientific assumptions \cite{exec24_framework_contracts}.

LangChain and LangGraph are retained as reference orchestration choices,
with conditional integration through typed causal-engine tools. Direct
adoption as the scientific product is deferred. Porting a small source
mechanism is an alternative only when its dependency and notice obligations
are understood and its semantics can be maintained. A custom controller
may be simpler for a narrowly bounded workflow; the comparison should
measure that alternative on the same recovery and refusal fixtures
\cite{exec24_framework_contracts}.

\subsection{AutoGen: conversational state and scientific state are different}
The retained AutoGen AgentChat 0.7.5 source says that its assistant retains
state across calls and should receive new messages rather than the whole
conversation again. It warns against concurrent calls on the same agent
instance. The documented tool behavior can involve concurrent tool calls
and an additional model call for reflection. Its save/load methods retain
model context; they do not establish transactional recovery of arbitrary
external side effects
\cite{v12_la04_autogen}.

For the retention example, separate conversational roles can propose
assumptions, challenge them and explain results. Agreement among those
roles is not independent scientific evidence if they share the same
model, data, prompt history or mistaken business premise. Their graph
proposals should enter a comparison set with provenance, not become
accepted orientations because a conversation reached consensus. Scientific
state must include the graph object, query, assumptions, identification
result and uncertainty object independently of the natural-language history
\cite{exec24_framework_contracts}.

The pinned README declares maintenance mode and recommends a successor.
That declaration informs a maintenance decision; the successor's advertised
properties are not established here by repeating the README. Reference
use and an existing-installation adapter remain possible, while a new
deployment decision needs an independently reviewed migration target and
matched fixtures. The proposal therefore defers unconditional adoption,
retains a conditional integration route, and records maintenance and
component-specific rights as explicit revisit conditions
\cite{v12_la04_autogen,exec24_framework_contracts}.

\subsection{Semantic Kernel: function filters do not verify scientific claims}
The inspected Python function-choice behavior exposes included/excluded
plugins and functions, and a maximum automatic-invocation setting. Its
kernel checks required and unexpected argument names, constructs an
invocation-filter stack, and can return a termination context. These are
useful controls for exposing only the intended causal tools at a given
stage. They do not verify the truth of a proposed exclusion restriction
or the adequacy of an adjustment set
\cite{exec24_framework_semantic_kernel}.

The selected automatic-invocation handler catches ordinary function
exceptions and places error text in a function result. A separate guard
is therefore needed before a later assistant message is accepted as a
qualified answer. The proposed protocol distinguishes an invalid request,
an unavailable service, a scientific refusal, a partial numerical result
and a successful calculation with stated assumptions. Argument-schema
validity is one condition among these, not their replacement
\cite{exec24_framework_semantic_kernel,exec24_framework_contracts}.

This review covers the selected Python profile. It does not transfer its
behavior to .NET, Java or every provider connector. Conditional integration
would expose a narrow tool set, retain scientific output objects and
test filters, failure results and termination using public fixtures.
Reference use is retained; broad cross-language parity and direct adoption
as the causal library are deferred. Porting requires separately inspecting
the actual component and maintaining its conversion contract
\cite{exec24_framework_contracts}.

\subsection{DSPy: optimizing a program requires a defensible evaluator}
The inspected ReAct implementation builds a trajectory, calls selected
tools and records execution exceptions as observations. It may leave the
tool loop after an invalid selection or context-window problem and then
run an extraction step. This is a useful mechanism for explaining partial
work, but the extraction step is not proof that the requested causal
calculation succeeded. The result must retain whether identification,
estimation and uncertainty checks actually completed
\cite{exec24_framework_dspy}.

The evaluation implementation supplies a particularly relevant positive
contract: after a returned evaluation run, missing per-example results
receive the configured failure score, and the aggregate denominator is
the full supplied development set. That guards against silently reporting
only successful examples in that code path. If the executor raises at
its error threshold, however, no completed aggregate should be invented.
The task definition, scorer, allowed failure score and failure threshold
must accompany every reported benchmark
\cite{exec24_framework_dspy}.

In an authored example, 1,000 requests produce 800 returned answers and
600 scientifically qualified answers. The qualified fraction among
returns is $600/800=3/4$, while the fraction among all requests is
$600/1000=3/5$. These answer different questions. Neither fraction alone
measures causal calibration, utility or the severity of unqualified
answers. A public experiment should report the complete request ledger,
refusals, errors, qualified answers and the prespecified loss assigned to
each outcome \cite{exec24_framework_contracts}.

The inspected parallel executor also deserves careful source-level
qualification. It can resubmit stragglers and stores completed outcomes
by input index, while a resubmission set is keyed by Future objects.
The nearby comment describing one resubmission per item is therefore not
by itself a proved per-item invariant: distinct Future identities can
refer to the same input. No native liveness or duplicate-work failure is
asserted here. A public test must trace the input key and every submission,
including work that continues after a useful result has been collected
\cite{exec24_framework_dspy,exec24_framework_contracts}.

DSPy is retained as a reference for program construction and evaluation,
with conditional integration around a fixed causal-task scorer. The
scorer must distinguish graph prediction, identification, numerical
estimation, uncertainty and business loss. Optimizing on the same examples
later used to claim performance would change the evaluation question.
Adoption as a causal oracle is rejected; bounded program optimization and
selective source reuse remain testable alternatives
\cite{exec24_framework_contracts}.

\subsection{LlamaIndex: stopping, structured output and runtime delegation}
The selected LlamaIndex base agent keeps iteration counts in context.
At its limit, the inspected path either raises or performs a final model
generation, depending on the chosen early-stopping method. A generated
summary after exhausting the budget is useful communication, but its
scientific status must remain partial or unresolved when necessary.
The selected function-agent code also distinguishes ordinary tool output
from a return-direct route and from handoff behavior
\cite{exec24_framework_llama_index}.

The separate workflows package makes the time contract especially clear.
Its selected constructor documents timeout in known-alive seconds accrued
through processed journal ticks; downtime during resumption and some
pre-crash time do not count. This is a different clock from an external
bank service deadline. In an authored example, two 20-second active
periods separated by a one-hour outage total 40 active seconds but
3,640 elapsed seconds. A caller needing a 60-second wall deadline cannot
infer it from the active-time limit
\cite{exec24_framework_workflows,exec24_framework_contracts}.

The inspected cancellation method signals the runtime and waits for the
result, while catching timeout, cancellation and other exceptions. Its
return therefore does not independently prove that every remote side
effect has stopped. The retry-policy source separately defines retry
predicates, delays and stopping conditions. These concrete distinctions
support an adapter design with external job reconciliation, explicit
deadline clocks and durable attempt identifiers. They do not establish
that the complete distributed runtime has been qualified
\cite{exec24_framework_workflows,exec24_framework_contracts}.

Conditional integration is attractive when document retrieval and a
workflow context are needed together, provided retrieved statements stay
traceable to inspected source material. Reference use is retained and
direct adoption as the scientific layer is deferred. The actual delegated
runtime, serializer, persistence backend and model connector need their
own fixture results. A source-level range match is only the start of
that engineering comparison \cite{exec24_framework_contracts}.

\subsection{Haystack: preserve explicit stop reasons and hook effects}
The selected Haystack agent supports state mappings, ordered hooks,
tool-concurrency control and a maximum number of agent steps. Its
documented step consists of a generator call and the requested tool work.
The inspected limit path records a maximum-steps exit reason and returns
the current public state. A tool-based exit checks error information;
the selected logic does not treat an errored exit-condition tool as a
successful exit. These are valuable details to preserve when comparing
returned results
\cite{exec24_framework_haystack}.

Hooks are scientifically consequential. A before-tool hook can alter
the pending request; an after-tool hook can transform returned messages;
an after-run hook can alter the final state. These facilities can reduce
context size or enforce a protocol, but a summary must not silently
discard an unresolved edge, confidence set or refusal. The selected
general function-based Tool class also lacks a formal output-name schema
by default, whereas specialized components can provide one. The proposed
adapter must therefore define and validate the causal result itself
\cite{exec24_framework_haystack,exec24_framework_contracts}.

For an authored example, the uncertainty sets $[-1,1]$ and $[-2,2]$
both have midpoint zero. A hook that keeps only the midpoint makes them
indistinguishable and cannot reconstruct either set from that scalar.
Whether a bank action is acceptable can depend on the endpoints even
when the midpoint is identical. Typed results should preserve the
original set, units, assumptions and numerical diagnostics; prose can
refer to that object without replacing it
\cite{exec24_framework_contracts}.

Haystack remains a reference and conditional pipeline choice. The
release-candidate source identity, tool schemas, hooks, state mappings
and runtime dependencies require matched public tests. Blanket adoption
as a causal identification engine is rejected. Selective reuse or a
smaller custom pipeline can be compared on the same preservation and
failure criteria \cite{exec24_framework_contracts}.

\subsection{A common public qualification design}
The comparisons suggest a shared test design rather than six incompatible
demonstrations. Start with the same business brief, data dictionary,
candidate graph set, formal query and causal-engine oracle. Change one
framework or backend factor at a time where attribution is intended;
retain the full factorial design when interactions matter. Test both
successful and deliberately unsupported cases: hidden-confounding
ambiguity, off-support interventions, unknown tool outcomes, malformed
schemas, interrupted state, duplicated acknowledgements and expired
deadlines. A persuasive narrative is not the scoring oracle
\cite{exec24_framework_contracts}.

Concurrency requires a testable budget invariant. If two workers each
read an unreserved balance of 100 USD and each accepts a 70 USD task,
both local checks pass while the combined commitment is 140 USD.
An atomic reservation rule can prevent that particular overspend;
a language-model instruction to stay within budget cannot establish
atomicity. The fixture must specify the transaction and recovery model,
then test reservations, release and repeated requests under that model
\cite{exec24_framework_contracts}.

The common adapter contract should carry immutable task and graph
identities, intervention and target, source citations, assumptions,
information timestamps, scientific status, numerical output, uncertainty,
logical action keys, attempts, deadlines and complete costs. It should
preserve uncertainty and refuse unsupported conversions. A CLI and an MCP
interface can transport this same contract; transport success does not
change the scientific verdict. Integrations with other repositories should
use their declared interfaces and public fixtures, with changes confined
to the proposed library's adapter until separate work is authorized
\cite{exec24_framework_contracts}.

Scientific ownership belongs to the target and assumptions; statistical
ownership to evaluation and dependence; adapter ownership to types,
refusals and replay; platform ownership to exact environments, rights,
budgets and recovery. Under hypothetical expenditure of 24,000 USD,
1,000 submitted, 800 returned and 600 qualified requests cost 24, 30 and
40 USD respectively. Actual provider prices, latency, native reliability
and bank value remain unmeasured
\cite{exec24_framework_contracts}.

The bounded reference choices are now concrete enough to guide a public
prototype comparison. They do not close the hosted/local backend survey,
the remaining framework and connector delegates, or the frontier literature
census. No framework removes unobserved confounding or recovers discarded
information without additional assumptions. Package execution and independent validation remain to be performed
\cite{exec24_framework_contracts}.


% SOURCE docs/release/2026-10-29/source/docs/writing-repair/2026-10-29/src/beb1ec47514d95d1.tex CHAR 0
\section{Named model backends: a scientific contract above the generation interface}
\label{sec:exec24-backend-census}

A bank considering an assistant for causal study design must choose both the
model service and the execution environment through which it is used. A useful
comparison asks what each environment preserves when a question becomes a
diagram, a tool request and a reviewed answer. It also asks what happens when
the output is partial, a tool fails, the request is retried or the budget is
exhausted. An API-shaped response is only one intermediate artifact. The
scientific object remains the question, information clock, population,
intervention, assumptions, graph and justified inference status.
\citep{exec24_backend_reference}

The selected primary profiles are the official Anthropic Python SDK 1.11.0,
Google GenAI Python SDK 2.28.0, the Ollama Python client, vLLM and llama.cpp.
Each repository revision is pinned no later than 7 October 2026. Ollama's
package version is derived from version-control metadata; the inspected
vLLM tree likewise does not establish a release wheel version. The llama.cpp
build declares base version 0.6.0 with development mode enabled by default.
These distinctions prevent a commit, a package release, a model name and a
deployed service from being treated as the same identity.
\citep{exec24_backend_anthropic,exec24_backend_google,exec24_backend_ollama,
exec24_backend_vllm,exec24_backend_llamacpp}

These five source profiles are deliberately bounded. Official OpenAI live-guide
requests returned HTTP 403; that omission does not establish the absence of a
service feature. The separate SDK and cookbook originals were subsequently
recovered and are examined in Section~\ref{sec:exec28-oa}. They resolve their
own source contracts without resolving the live-guide omission. Hosted-service
entitlements, immutable service revisions, complete model catalogues, regional
behavior and actual prices remain outside these source-profile conclusions. Azure and Bedrock deployments, additional
model hosts and other local engines require their own applicable evidence.
The following choices therefore supply a prototype route while leaving the
broader backend census open.~\citep{exec24_backend_reference}

\subsection{Separate syntax, object validation and scientific validity}

Let $x$ denote a generated representation and let $S(x)$ mean that it satisfies
a declared serialization schema. Let $G(x)$ mean that it denotes an admissible
graph and $I(x)$ mean that the intervention query is justified under the
declared causal assumptions. These are different predicates. A conventional
edge-list schema can admit the three edges $A\to B$, $B\to C$ and $C\to A$:
every endpoint is a valid named variable and every field has the expected
type, but the graph is cyclic. An independent topological check rejects it
when the contract requires a DAG. Even an acyclic graph does not establish
that it represents the bank's assignment process. The product must therefore
preserve separate schema, graph and identification decisions rather than
promoting $S(x)$ directly to $I(x)$.
\citep{exec24_backend_reference}

The Anthropic SDK illustrates another level of this distinction. Its selected
output configuration contains a JSON-schema format field. Separately, the
client's internal \texttt{\_strict\_response\_validation} option defaults to
false. The response-processing routine invokes a type validator when that
option is enabled and otherwise constructs the target object; a reported
validation error is wrapped as an API response-validation error. These are
client object-handling branches, not a causal validity oracle or proof about
every service-side generation path. A proposed adapter keeps the requested
schema, returned content and observed validation status separately.
\citep{exec24_backend_anthropic}

Google's selected response conversion is more explicit about the requested
Python representation. A Pydantic model invokes its JSON validator; caught
JSON or validation errors leave the parsed result unset. A dictionary or
SDK Schema branch instead applies JSON parsing to the returned text, without
an additional local full-schema validator in that branch. A caller must
not interpret a returned response object, or merely successful JSON parsing,
as proof that all requested schema constraints hold. These are inspected
client branches, not measurements of service violation rates.
\citep{exec24_backend_google}

For vLLM, the pinned documentation distinguishes schema-derived decoding
from structural-tag parsing. In the latter route, constraining the tool-call
envelope and enforcing an individual tool's argument schema are separate.
The documented per-tool strict flag, server strictness setting and global
enforcement switch affect that route. General named-function prose is
qualified by this later parser-specific explanation; it is unsafe to repeat
the broad statement as an unconditional guarantee across all parsers.
The structured-output configuration also chooses a backend automatically
by default and describes that choice as release-dependent. A prototype must
pin the parser and decoding backend as well as the model and server revision.
\citep{exec24_backend_vllm}

llama.cpp exposes grammar and JSON-schema generation controls and directs
readers to supported-feature tests. That is a useful constrained-generation
boundary, but neither the parameter name nor a syntactically valid output
establishes full JSON Schema coverage, a valid causal graph or identification.
The proposed common validator runs after every backend and retains the
backend's original output for diagnosis. Unsupported constraints must be
reported, not silently weakened to make a request run.
\citep{exec24_backend_llamacpp,exec24_backend_reference}

\subsection{Tool requests, tool execution and scientific responsibility}

The Google SDK documents automatic function calling when a Python function
is supplied directly. The inspected loop counts remote model calls, retrieves
function-call parts and invokes the registered functions through a helper.
At the last permitted remote call it returns the model's function calls
without executing them, because no remote call remains for their results.
The helper can visit several function-call parts in a single response and
can convert a local exception into an error-bearing function response.
Consequently a remote-call limit is not a limit on the number, duration or
economic consequences of local function invocations.
\citep{exec24_backend_google}

In an authored example, two model responses are permitted. The first contains
three executable tool calls; the final response contains two more that remain
unexecuted under the stipulated last-call rule. There are two remote calls,
three executed local functions and two unresolved requests. With hypothetical
model costs of 0.20 and 0.30 USD and tool costs of 2, 5 and 3 USD, expenditure
is 10.50 USD. Reporting only two calls loses both execution and cost meaning.
For a credit-policy assistant, the adapter should reserve separate budgets
for model attempts, scientific tools, human review, wall time and monetary
expenditure before allowing an action.
\citep{exec24_backend_reference}

The selected Ollama \texttt{chat} path accepts tools, converts their
declarations, sends the request and returns a typed response or a stream of
typed chunks. That path does not itself implement the Google-style local
function-execution loop. Applications using the two interfaces therefore
need different orchestration policies even when both accept a tools argument.
For the proposed library, generation may suggest a causal operation, while
the controller checks its scientific context and permission before execution.
A model request for estimation cannot create identification that the causal
engine has refused.~\citep{exec24_backend_ollama,exec24_backend_reference}

\subsection{Completion, refusal and observation deadlines}

The Anthropic message type distinguishes a natural end, token limit,
specified stop sequence, tool request, paused turn, refusal and context-window
limit. Its documentation says that a streaming message begins without a final
stop reason. Google's finish-reason type likewise separates ordinary stopping,
token limits, safety filtering, malformed function calls and other outcomes.
Ollama's response model carries \texttt{done}, \texttt{done\_reason}, token
counts and durations expressed in nanoseconds. A common adapter should
retain native states before mapping them to its own completion, partial,
refused or failed categories. Provider-policy refusal is distinct from an
engine's scientific refusal to identify an effect.
\citep{exec24_backend_anthropic,exec24_backend_google,exec24_backend_ollama}

Time budgets need equally explicit units and scope. The Anthropic constants
set a ten-minute default HTTP timeout profile and two retries; the inspected
retry routine considers selected HTTP errors and server retry hints. Google's
HTTP options describe timeout in milliseconds and maximum attempts including
the initial request. Ollama's client explicitly defaults timeout to none.
These observations do not establish a whole-business-workflow deadline.
Retries, model generation, local tools, queueing and result review require a
separate end-to-end clock and a declared policy for unresolved remote work.
\citep{exec24_backend_anthropic,exec24_backend_google,exec24_backend_ollama}

The distinction has exact consequences even without calling a provider. Two
retries permit three request attempts; three stipulated four-second attempts
and two twelve-second retry waits total 36 seconds. Calling this a
four-second operation because each observed request took four seconds would
omit most elapsed time. This example is accounting under a specified sequence,
not a measured latency or a universal upper bound.
\citep{exec24_backend_reference}

The selected llama.cpp source makes the prediction-time qualification
concrete. The inspected stopping branch tests elapsed generation time inside
the branch that sends a newline. Its documented limit concerns prediction,
not queueing and prompt processing. With hypothetical queueing of five
seconds, prefill of three seconds and the next relevant newline at eight
seconds of generation, a two-second generation setting can be observed at
sixteen seconds end to end, assuming no earlier stop. This is a conditional
control-flow witness, not a benchmark. The inspected Cancel task releases a
matching slot; that operation alone does not reverse a previously executed
external tool or establish every transport-disconnection behavior.
\citep{exec24_backend_llamacpp,exec24_backend_reference}

\subsection{Compatible endpoints and local deployment are separate choices}

vLLM documents several OpenAI-compatible routes while identifying parameter
exceptions and model-dependent tool behavior. Its server can apply sampling
defaults from a model repository's generation configuration unless configured
otherwise. llama.cpp exposes compatible routes and warns that prompt-cache
reuse can yield nonidentical logits across batch sizes on some backends.
Thus matching an endpoint and seed does not establish identical generation
semantics or bitwise reproducibility. A reproducibility record binds model
weights, tokenizer, chat template, quantization, decoding backend, sampling
configuration, cache and hardware conditions.
\citep{exec24_backend_vllm,exec24_backend_llamacpp}

The direct source notices differ from the rights of those artifacts.
Anthropic's Python SDK, the Ollama client and llama.cpp carry MIT notices;
Google GenAI and vLLM declare Apache 2.0. These observations concern selected
source distributions. They do not provide hosted-service access, model-weight
rights, training-data permissions, a redistributable dependency environment
or a bank support agreement. The proposed first comparison references these
implementations, conditionally integrates an original client or local worker
after qualification, and rejects an unmotivated wholesale port. Rewriting a
client would preserve the missing service and scientific evidence while adding
a second maintenance burden.
\citep{exec24_backend_anthropic,exec24_backend_google,exec24_backend_ollama,
exec24_backend_vllm,exec24_backend_llamacpp,exec24_backend_reference}

\subsection{A public comparison that can justify a prototype}

The first experiment should use public or generated business-to-causal
questions with explicit target, graph and refusal oracles. The same logical
cases, evidence and rubrics should be paired across backends; repeated attempts
measure within-case variation rather than creating new independent cases.
Record all submissions, completed responses, parseable objects, scientifically
acceptable answers, appropriate refusals and unresolved failures.
\citep{exec24_backend_reference}

For illustration, backend A correctly answers 45 of 50 completed cases from
100 submissions, while B correctly answers 72 of 90 completed cases from
the same number submitted. Conditional agreement favors A, 90 versus
80 percent; correct answers per submitted case favor B, 72 versus 45 percent.
At hypothetical total costs of 150 and 180 USD, cost per correct answer is
$10/3$ versus $5/2$ USD. These invented data show why the evaluation needs
both conditional quality and coverage; they rank no real provider.
\citep{exec24_backend_reference}

The scientific owner specifies the cases, loss and acceptance standard;
the adapter owner preserves state, units and artifacts; the platform owner
qualifies resource and rights conditions; and the evaluation owner accounts
for attempts, exclusions and full cost. Broader hosted evidence and native
trials remain necessary before a named backend is selected for institutional
use. The comparison makes the next public experiment concrete without
claiming that client code, constrained generation or a completed request
constitutes a qualified causal product.
\citep{exec24_backend_reference}


% SOURCE docs/release/2026-10-29/source/docs/writing-repair/2026-10-29/src/e0aa42e85efb2ce0.tex CHAR 0
\section[Structured outputs and scientific admission]{Interpreting structured answers from a language model}
\label{sec:exec28-oa}
A structured banking answer must preserve both the scientific target and the state of the computation.
Suppose a collections analyst requests the effect of a reminder on thirty-day arrears. A response can have
the required fields and a well-formed number while using a post-treatment variable, an unjustified graph,
or a different treatment population. It can also be a syntactically valid fragment of an unfinished request.
The proposed library therefore needs separate tests for completion, representation and causal admission.
This distinction determines how a provider client can be reused without attributing scientific guarantees to it.

\subsection{What the pinned official sources establish}
The selected official OpenAI cookbook demonstrates a strict JSON-schema response format, required fields,
closed object properties and a Pydantic parsing helper. Its worked examples are useful interface illustrations,
not measurements of this proposal. The same notebook explicitly describes a refusal field because a refused
request need not follow the supplied answer schema. Its opening description of constrained outputs must
therefore be read together with this exception. Schema adherence concerns representation; it does not assert
that an adjustment set is scientifically admissible or that a reported effect is identified.~\citep{exec28_oa_cookbook}

The selected Python SDK source has version 3.26.0 and declares Python 3.10 or later. Its parsing helper converts
a Pydantic model to a schema and parses eligible content into that model. In the inspected Chat Completions path,
a length stop raises a dedicated error, as does a content-filter stop. Refusal content is not parsed as an ordinary
answer. Automatic tool-argument parsing requires a strict function-tool contract. These are source-level client
rules at the selected revision, rather than a claim that every compatible server or model honors them.
The cookbook and SDK are pinned separately; their example calling surfaces must not be silently combined into
one version-independent executable recipe.~\citep{exec28_oa_sdk,exec28_oa_cookbook}

The Responses parsing path also makes the message phase significant. At this revision, final-answer text and
legacy messages without a phase are eligible for parsing; other explicit phases retain their text with no parsed
result. The helper documentation warns that concatenated output text can include commentary, while the parsed
output selects a structured result. An adapter should retain the entire typed output sequence and the response
status, not search arbitrary text for a schema-shaped substring and promote it to the final answer. A declared
incomplete state remains incomplete even if part of its output can be decoded.~\citep{exec28_oa_sdk}

For the banking request, a useful envelope contains the provider response identifier, endpoint and selected
model identifier; status and native stop reason; refusal or incomplete details; the parsed payload, if any;
tool-call identities and pending results; and the target, assumptions and scientific disposition. These fields
are proposed product requirements. A provider refusal, a malformed object and a correctly identified failure
of identification have different meanings. The analyst should be able to tell which information is missing
and which action, if any, could change the result.

\subsection{An admission rule that cannot be satisfied by JSON alone}
Let \(C\) denote completed output, \(J\) successful representation validation and \(I\) admission of the
requested causal effect under its declared information and assumptions. A proposed effect-admission predicate is
\[
 A=C\land J\land I.
\]
This is a software decision rule, not a theorem that the assumptions are true. It makes the owner of each
judgment visible. A source parser can help establish \(J\), transport and response-state handling contribute to
\(C\), and the causal library must justify \(I\). The operator and eventual bank reviewer still own the truth
of substantive inputs. An unidentified effect may legitimately produce a completed, well-formed scientific
answer whose disposition is nonidentification; it must not produce an admitted scalar effect by default.

An exact authored four-case witness admits only the complete, schema-valid, identified case. A complete but
unidentified case retains its nonidentification result; an incomplete case with a valid prefix remains pending;
a refusal stays a refusal. These cases are finite constructions used to specify the interface, not observations
of OpenAI behavior or a measured refusal rate. In particular, treating all valid objects as completed causal
answers would merge scientifically different outcomes.~\citep{exec28_oa_reference}

\subsection{Retry layers and the cost of one business request}
The pinned SDK documentation describes two retries by default for selected connection and HTTP failures,
subject to a resendable body, and a configurable request timeout. The implementation distinguishes status
classes, server retry hints and backoff. A workflow can add another retry loop around this client. Consequently,
neither the number of logical banking requests nor one configured timeout determines the full attempt count or
end-to-end latency. A retry may also leave uncertainty about a remote effect; successful parsing alone cannot
provide exactly-once execution of a business action.~\citep{exec28_oa_sdk}

Consider a stipulated workflow with at most two outer attempts and three HTTP attempts within each. Its maximum
is six HTTP attempts, not the sum of the two limits. At a stipulated USD 0.25 per attempt, the corresponding
cost is USD 1.50. If each attempt takes four seconds, inner waits are one and two seconds in each outer attempt,
and the outer wait is three seconds, the specified path takes \(24+6+3=33\) seconds. These values are exact
bookkeeping under the stated assumptions, not vendor prices or a measured service guarantee.~\citep{exec28_oa_reference}

The proposed admission layer should reserve a workflow-wide attempt and cost budget before issuing calls,
debit actual attempts, retain response identifiers and unknown outcomes, and distinguish cancellation requested
from cancellation observed. Public prototype tests should include an error after remote work may have begun,
an incomplete stream, a refusal, invalid tool arguments and a timeout during a tool result. The complete
denominator includes these cases. Financial value remains a separate assessment of decision quality and full
workflow costs; cheap valid JSON is not a monetary benefit by definition.

\subsection{Comparison and a bounded reuse choice}
The existing comparison treats Anthropic and Google client contracts, Ollama request and stream behavior,
and the parser, template and generation settings of vLLM and llama.cpp separately. The newly inspected OpenAI
originals add a sixth named provider/client family with explicit parsing, refusal, phase and retry rules.
Across these profiles, one adapter should preserve native status and error information before mapping it into
the common scientific envelope. Compatibility in endpoint spelling cannot prove equivalence in parsing,
model behavior, cancellation or cost.~\citep{exec24_backend_anthropic,exec24_backend_google,exec24_backend_ollama,exec24_backend_vllm,exec24_backend_llamacpp,exec28_oa_sdk}

The bounded proposal choice is conditional adoption of the original client behind a thin, versioned adapter,
rather than a wholesale port. The inspected SDK and cookbook carry Apache 2.0 and MIT notices respectively;
these notices concern those source artifacts. They do not establish service entitlement, model-weight rights,
data permissions or approved bank deployment. An implementation decision still requires the exact endpoint,
model, dependencies, configured retry policy and applicable service terms. The official live guide access
failure also remains an omission; the recovered repository originals establish their own narrower contracts.
No downloaded code or live model was executed in this review.~\citep{exec28_oa_sdk,exec28_oa_cookbook}

CLI and MCP should expose the same typed outcome, including nonidentification, scientific rejection, refusal,
incomplete output and transport uncertainty. Framework scheduling must consume this outcome rather than infer
success from text presence. This supplies a concrete route for a low-cost public prototype: freeze the client
and adapter identities, register adversarial response-state fixtures, verify the scientific envelope against
independent answers, and only then compare live backends under a declared budget. The source comparison reduces
an architectural uncertainty while leaving native performance, service terms and bank effectiveness open.
The banking decision remains admissible only when its scientific meaning survives the complete request lifecycle.


% SOURCE docs/release/2026-10-29/source/docs/writing-repair/2026-10-29/src/cf38655a502a0afa.tex CHAR 0
\section[Framework delegates and scientific completion]{When a completed software call leaves the research unfinished}
\label{sec:exec28-dl}
A causal workflow can finish its control flow without producing an admissible causal answer.
For a collections analyst, the distinction is practical. A reminder-policy analysis may return an
object after its structured-output step failed, repeat a tool after an interrupted response, or
accept a treatment identifier under a weakened constraint. Each event changes what the analyst can
infer from the result. A framework comparison must therefore follow the delegated implementation
of validation, stopping, retry and recovery, rather than infer these properties from a common API.
The following comparison concerns selected, pinned original source paths. It motivates product
contracts and public qualification experiments; it does not report executed vendor integrations.

\subsection{Validation must survive the complete calling path}
LangChain's inspected tool base distinguishes a dictionary schema from a Pydantic model. In the
dictionary branch of the input parser, the supplied input is returned; the Pydantic branch invokes
model validation and handles an injected tool-call identifier. This local branch difference does
not prove that no other layer validates a dictionary-schema request. It does show why the proposed
adapter cannot treat the presence of a schema as evidence that validation occurred at every entry
point. The same source can convert a handled validation or tool exception into an error-status tool
message. The consumer must preserve that status even when the returned content is readable text.
\citep{exec28_dl_langchain}

DSPy's inspected tool helper iterates over provided arguments, conditionally validates a field
against its schema when a qualifying type is present, then uses typed parsing for applicable
arguments. Its synchronous and asynchronous call paths also differ in coroutine handling. Google
GenAI's function-declaration helper builds a schema from a callable signature and type information,
including required fields and definitions; an inspected Vertex conversion branch supplies an
object type for an otherwise untyped union. These are useful mechanisms for translating interfaces.
They do not establish that a variable precedes treatment, that a graph encodes the business question,
or that omitted information is scientifically harmless. Conditional adoption should retain the
original helper while placing a common causal-request validator at the library boundary.
\citep{exec28_dl_dspy,exec28_dl_google}

Consider a request whose treatment identifier must contain a digit. Within the stipulated domain
\(D=\{\texttt{x},\texttt{7}\}\), the intended accepted set is
\(S=\{\texttt{7}\}\). If a conversion relaxes the restriction to any string, the accepted set becomes
\(\widetilde S=D\), so \(S\subsetneq\widetilde S\). The extra value is a false admission under the
original contract. This finite authored witness establishes the logical consequence of weakening;
it is not an observed generation or a measured package failure.~\citep{exec28_dl_reference}

The distinction matters for llama.cpp's inspected schema-to-grammar converter. An unsupported
regular-expression pattern can generate a warning and fall back to any string, whereas an invalid
pattern is recorded as an error. The converter subsequently raises recorded errors and prints
warnings. A separate build/configuration path can use LLGuidance; the warning behavior must not be
generalized to every backend. The selected chat-template path chooses schema-derived grammar or a
supplied grammar, while a separate argument parser can retain text when JSON parsing fails.
Consequently, grammar generation, parse success and validation against the original schema are
three distinct observations. The proposed library should retain conversion diagnostics and the
selected grammar backend, and validate the completed payload against the original request before
admitting an effect. Wholesale porting would inherit these semantic choices and add maintenance
cost without removing the need for this validation.~\citep{exec28_dl_llamacpp}

\subsection{A stop event is a control state, not a scientific certificate}
In LlamaIndex's inspected agent workflow, structured-response generation can raise an exception
that is converted to a warning before the workflow returns a stop event. Its early-stop generation
option can also make another model call to construct a final response. Thus, a terminal event alone
cannot certify a valid structured answer. In LangChain's classic agent, forced stopping returns a
limit response, while the generation option can request a final model answer. The synchronous loop
updates elapsed time between steps; that observation is narrower than a claim that all outstanding
remote operations are cancelled at a wall-clock limit.~\citep{exec28_dl_llama_index,exec28_dl_langchain}

Let \(T\) indicate terminal control flow, \(V\) validation of the scientific envelope, and \(I\)
admission of the requested effect under declared assumptions. The proposed scalar-effect gate is
\(T\land V\land I\). An exact four-case construction admits only the terminal, valid, identified
case. A completed nonidentification result remains a legitimate answer with a different disposition;
a terminal structured-output failure remains a failed representation; an incomplete stream remains
incomplete. The gate specifies responsibility, not the truth of a bank's assumptions. A framework
owns its execution state, the adapter checks representation and preserves errors, and the causal
library supplies the identification argument.~\citep{exec28_dl_reference}

Semantic Kernel's inspected chat agent combines arguments and execution settings before delegating
to a selected provider. Its function wrapper invokes filters, records errors and propagates
exceptions; the streaming path can yield partial results. These hooks are useful integration
locations, but the inspected wrapper does not establish a causal validator. Haystack validates
components, sockets and mandatory inputs; its execution wrapper checks the returned mapping and
output names, and a blocked queue can raise an error. Such checks protect pipeline structure.
They should be retained alongside a separate validation of the target, information clock, graph,
estimand and uncertainty claim. A proposed CLI command and MCP tool should return the same typed
scientific disposition even when their orchestration frameworks differ.
\citep{exec28_dl_semantic_kernel,exec28_dl_haystack}

\subsection{Budgets must count actual attempts and side effects}
LangChain recommends placing retries around the smallest operation that needs them. Its inspected
retry wrapper has configurable exception handling and backoff; the batch path retains successful
items and retries failed original indices. Suppose three independent items \(a,b,c\) run once,
only \(b\) fails, and \(b\) needs two further attempts. Selective retry makes \(1+3+1=5\) item
attempts, whereas repeating the whole three-item batch three times makes nine. At a stipulated
USD 0.25 per item attempt, the costs are USD 1.25 and USD 2.25. These exact accounting examples
are not vendor prices. They also assume that retained successes are valid and that the repeated
operation is safe to retry; neither property follows merely from an exception class.
\citep{exec28_dl_langchain,exec28_dl_reference}

Google GenAI's inspected HTTP path converts configured milliseconds into seconds and supports
request-specific retry settings in addition to client defaults. The inspected server-timeout header
rounds seconds upward. Haystack's visit limit applies to an individual component, while a workflow
can traverse several components. A business-request budget must therefore count HTTP attempts,
tool attempts, generation calls, waiting time and unresolved outcomes explicitly. A framework
iteration cap, a component visit cap and an HTTP timeout are different limits.
\citep{exec28_dl_google,exec28_dl_haystack}

Anthropic's pinned tool-runner documentation describes dispatch by native stop reason. Tool-use
responses trigger tools; paused turns and compaction can require continuation; other stop states
terminate the iterator after yielding the message. Its documented eager streaming mode can execute
a tool before the overall reply finishes. If that reply is interrupted, the tool's effect cannot be
undone merely by discarding the conversation, and the model may request the call again. The same
documentation describes compaction calls outside the ordinary iteration count and nested usage
accounting. Two ordinary turns plus one separate compaction make three calls in an authored counting
example. A top-level counter cannot be assumed to represent the complete cost denominator.
\citep{exec28_dl_anthropic,exec28_dl_reference}

For the reminder-policy workflow, graph inspection and effect calculation can be read-only tools;
dispatching an actual reminder is a different operation. The proposal should assign stable operation
identities, preserve pending results, and require a deduplication or reconciliation policy before
retrying actions with external effects. Cancellation requested and cancellation observed must be
separate states. A failed or unknown outcome remains in the denominator when comparing cost and
reliability. These are proposed architectural requirements, not a claim that any reviewed framework
already supplies end-to-end exactly-once execution.

\subsection{Recovery, native stop reasons and learned programs}
The inspected delegated workflow context describes a process-local store and serialization of
state, event queues, buffers and related execution metadata. Cancellation is delegated for an
external context, while other context states issue warnings. Serialization can support resumption,
but this source path alone cannot establish that a remote action and a local checkpoint commit
atomically. The proposed recovery record must therefore include both local state and unresolved
external-operation identities. Tests should interrupt immediately before and after issuing a tool,
receiving its result and recording the result.~\citep{exec28_dl_workflows}

Native stop reasons also need to survive compatibility adapters. Anthropic's selected type includes
refusal, pause, tool use, output limits and context exhaustion. In vLLM's inspected serving path,
parser selection depends on tool and reasoning settings, strictness, the model and tokenizer.
Its compatibility finish reason can depend on tool mode and generated tool content while retaining
the native stop reason. llama.cpp distinguishes native end-of-sequence, stop-word and limit states
before constructing compatibility fields; the task structure also retains cancellation targets and
parent-child identities. These source representations support a concrete preservation requirement.
They do not independently prove that cancellation has reached every child or that two servers with
similar response fields provide equivalent behavior.~\citep{exec28_dl_anthropic,exec28_dl_vllm,exec28_dl_llamacpp}

Optimization introduces another boundary. DSPy's inspected bootstrap procedure selects traces from
training examples. A variable called validation contains the unbootstrapped training subset, and
success depends on the configured metric; without a metric, the inspected branch accepts success.
The threshold branch also depends on its truth value. These are selected source semantics, not a
claim of a general empirical defect. A causal assistant still needs separately held-out tasks with
independent expected answers, including nonidentification, invalid adjustment and ambiguous graph
cases. Optimizing a demonstration-selection metric cannot by itself establish causal validity on
new banking questions.~\citep{exec28_dl_dspy}

\subsection{A bounded adoption and qualification decision}
The proposal favors thin adapters around selected original clients and orchestration components.
LangChain contributes composable calls and scoped retries; Semantic Kernel contributes provider and
function hooks; Haystack contributes explicit pipeline structure; LlamaIndex and its delegated
workflow package contribute orchestration and recovery representations; DSPy contributes a candidate
optimization layer subject to an independent evaluation contract. The decision is conditional
integration, with scientific validation owned by the causal library. It rejects automatic promotion
of terminal messages into effects, implicit schema weakening, unmeasured retry economics and
unqualified wholesale ports.~\citep{exec28_dl_langchain,exec28_dl_semantic_kernel,exec28_dl_haystack,exec28_dl_llama_index,exec28_dl_workflows,exec28_dl_dspy}

For hosted backends, client source inspection and service qualification remain separate. For local
backends, the exact server build, model, tokenizer, template, parser and grammar configuration are
part of the candidate. Ollama's selected README itself requires a separately running server and a
selected model; its client identity cannot substitute for those artifacts. README migration or
product-focus statements likewise supply leads, not inherited qualification: the Semantic Kernel
README names a successor framework, and the LlamaIndex README names parsing products and benchmarks.
Those separate products require their own primary-source comparison before adoption claims.
\citep{exec28_dl_ollama,exec28_dl_semantic_kernel,exec28_dl_llama_index}

A low-cost public qualification suite can now be specified precisely. It should present an invalid
original-schema value, a handled tool error, a structured-output failure followed by a terminal event,
a partly failed batch, an interrupted side-effecting call, a compaction call outside the ordinary
counter, and a limit stop containing a valid-looking prefix. Each fixture needs an independent expected
scientific disposition, attempt ledger and recovery outcome. Subsequent native trials require exact
artifacts, rights and declared budgets; bank effectiveness remains a later experiment. The comparison
supports these targeted experiments without converting static source inspection into measured merit.
The causal answer is trustworthy only when its meaning survives validation, scheduling, stopping and recovery.

\subsection{The component choice follows its scientific contract}

The complete comparison makes a conditional architecture possible without
assuming that every attractive component belongs in it. Typed representations
can preserve an effect set; source and state identities can expose stale work;
a deterministic controller can retain refusals and unknown outcomes. The
evaluation, memory, specialist-interface and code-generation studies show why
these properties still need task-specific tests, independent references and
complete cost accounting. Their distinct failure mechanisms determine what
should be checked before reuse.

The initial investment is therefore a bounded comparison against the existing
controller and a structured analyst workflow. A selected original client or
framework needs its own exact environment and behavioral qualification. No
source profile supplies causal truth, a universal performance ranking or bank
acceptance. The following synthesis draws together the assistance and
implementation choices while retaining those limits.


